\documentclass[aspectratio=169,11pt]{beamer}

\usetheme{Madrid}
\usecolortheme{seahorse}
\setbeamertemplate{navigation symbols}{}
\setbeamertemplate{footline}[frame number]

\usepackage[utf8]{inputenc}
\usepackage[T1]{fontenc}
\usepackage[spanish,es-tabla]{babel}
\usepackage{booktabs}
\usepackage{array}
\usepackage{amssymb}
\usepackage{amsmath}
\usepackage{tikz}
\usepackage{pgfplots}
\pgfplotsset{compat=1.18}
\usetikzlibrary{arrows.meta,positioning,shapes.geometric}

\title{Guía de estudio: Estadística, Probabilidades y Modelos Estocásticos}
\subtitle{Descriptiva, índices, probabilidad, variables aleatorias y dispersión probabilística}
\author{Material de apoyo para estudio}
\date{\today}

\newcommand{\pregunta}[1]{\textbf{Pregunta que responde:} \emph{#1}}
\newcommand{\errorcomun}[1]{\textbf{Error común:} #1}
\newcommand{\diferencia}[1]{\textbf{Diferencia clave:} #1}
\newcommand{\intuicion}[1]{\textbf{Intuición:} #1}
\newcommand{\observacion}[1]{\begin{block}{Observación}#1\end{block}}

\begin{document}

\begin{frame}
  \titlepage
\end{frame}

\begin{frame}{Cómo usar esta guía}
\begin{itemize}
  \item Cada diapositiva aclara un concepto central de estadística descriptiva.
  \item La estructura se repite: definición formal, intuición, pregunta, ejemplo, diferencias y errores comunes.
  \item El foco está en conceptos parecidos que suelen confundirse en evaluaciones y ejercicios.
  \item La guía integra estadística descriptiva, números índices, teoría de probabilidades y variables aleatorias.
\end{itemize}
\end{frame}

\begin{frame}{Mapa general de tópicos}
\begin{columns}[T]
\column{0.48\textwidth}
\begin{enumerate}
  \item Introducción a la estadística
  \item Conceptos fundamentales
  \item Método estadístico
  \item Clasificación de variables
  \item Escalas de medición
\end{enumerate}
\column{0.48\textwidth}
\begin{enumerate}\setcounter{enumi}{5}
  \item Dimensionalidad
  \item Estadística descriptiva
  \item Tablas de frecuencia
  \item Variables continuas e intervalos
  \item Tablas de contingencia
  \item Medidas de posición y dispersión
  \item Números índices económicos
  \item Probabilidades y análisis combinatorio
  \item Variables aleatorias y modelos estocásticos
\end{enumerate}
\end{columns}
\end{frame}

\section{Introducción}

\begin{frame}{¿Qué es la estadística?}
\textbf{Definición formal:} disciplina que recolecta, organiza, resume, analiza e interpreta datos para apoyar conclusiones y decisiones.
\vspace{0.4em}
\begin{itemize}
  \item \textbf{Intuición:} convierte datos sueltos en evidencia comprensible.
  \item \pregunta{¿Qué dicen los datos sobre el fenómeno que estudio?}
  \item \textbf{Ejemplo:} resumir las notas de un curso mediante promedio, mediana y gráficos.
  \item \errorcomun{creer que la estadística solo consiste en hacer cálculos; también implica diseñar, interpretar y comunicar.}
\end{itemize}
\end{frame}

\begin{frame}{Estadística y toma de decisiones}
\textbf{Idea central:} la estadística reduce incertidumbre, pero no elimina totalmente el riesgo.
\begin{block}{Flujo de decisión}
Datos $\rightarrow$ análisis estadístico $\rightarrow$ información $\rightarrow$ decisión
\end{block}
\begin{itemize}
  \item \textbf{Ejemplo:} una universidad analiza tasas de aprobación para decidir si refuerza un curso.
  \item \diferencia{una decisión estadísticamente informada usa evidencia; una decisión intuitiva depende principalmente de percepción o experiencia.}
  \item \errorcomun{confundir una conclusión probable con una certeza absoluta.}
\end{itemize}
\end{frame}

\begin{frame}{Estadística y método científico}
\begin{enumerate}
  \item Plantear problema.
  \item Formular hipótesis.
  \item Diseñar el estudio.
  \item Recolectar datos.
  \item Analizar resultados.
  \item Concluir y comunicar.
\end{enumerate}
\begin{itemize}
  \item \textbf{Rol de la estadística:} conecta observaciones con conclusiones verificables.
  \item \pregunta{¿La evidencia empírica apoya la hipótesis?}
  \item \errorcomun{analizar datos sin definir antes el problema o la hipótesis.}
\end{itemize}
\end{frame}

\begin{frame}{Descriptiva vs. inferencia estadística}
\begin{columns}[T]
\column{0.48\textwidth}
\textbf{Estadística descriptiva}
\begin{itemize}
  \item Organiza y resume datos observados.
  \item Usa tablas, gráficos y estadísticos.
  \item No generaliza necesariamente a datos no observados.
\end{itemize}
\column{0.48\textwidth}
\textbf{Inferencia estadística}
\begin{itemize}
  \item Usa muestras para concluir sobre poblaciones.
  \item Incluye estimación y pruebas de hipótesis.
  \item Trabaja explícitamente con incertidumbre.
\end{itemize}
\end{columns}
\vspace{0.4em}
\diferencia{describir es resumir lo observado; inferir es proyectar hacia lo no observado.}
\end{frame}

\section{Conceptos fundamentales}

\begin{frame}{Universo}
\textbf{Definición formal:} conjunto de individuos, unidades o elementos que forman el objeto de estudio.
\begin{itemize}
  \item \textbf{Intuición:} las personas, objetos o casos sobre los que se quiere estudiar algo.
  \item \pregunta{¿Quiénes forman parte del estudio?}
  \item \textbf{Ejemplo:} todos los alumnos de la UTFSM.
  \item \diferencia{el universo contiene individuos; la población contiene valores de una variable medida sobre esos individuos.}
  \item \errorcomun{usar ``universo'' y ``población'' como sinónimos sin distinguir individuos y valores.}
\end{itemize}
\end{frame}

\begin{frame}{Variable}
\textbf{Definición formal:} característica, atributo o magnitud que se mide u observa en los individuos del universo.
\begin{itemize}
  \item \textbf{Intuición:} lo que se quiere conocer de cada individuo.
  \item \pregunta{¿Qué estoy midiendo u observando?}
  \item \textbf{Ejemplos:} edad, estatura, carrera, promedio académico.
  \item \diferencia{la variable es la característica; el dato es un valor particular de esa característica.}
  \item \errorcomun{confundir ``edad'' con ``20 años'': edad es variable; 20 es dato.}
\end{itemize}
\end{frame}

\begin{frame}{Población}
\textbf{Definición formal:} conjunto de todos los valores que toma una variable específica en todos los individuos del universo.
\begin{itemize}
  \item \textbf{Intuición:} la lista completa de respuestas posibles observadas para una variable en todos los individuos relevantes.
  \item \pregunta{¿Cuáles son todos los valores de la variable en el universo?}
  \item \textbf{Ejemplo:} si la variable es edad, la población es \{20, 22, 19, \dots\} para todos los alumnos.
  \item \diferencia{universo: alumnos; variable: edad; población: edades de todos los alumnos.}
  \item \errorcomun{decir que la población son ``los alumnos'' cuando, en esta distinción, son los valores medidos.}
\end{itemize}
\end{frame}

\begin{frame}{Muestra}
\textbf{Definición formal:} subconjunto de la población que se observa o mide efectivamente.
\begin{itemize}
  \item \textbf{Intuición:} una parte representativa, o al menos seleccionada, de los datos completos.
  \item \pregunta{¿A quiénes o qué valores medí realmente?}
  \item \textbf{Ejemplo:} edades de 100 alumnos seleccionados de la UTFSM.
  \item \diferencia{la población es el conjunto completo de valores; la muestra es solo una parte.}
  \item \errorcomun{asumir que cualquier muestra representa bien a la población sin revisar cómo fue seleccionada.}
\end{itemize}
\end{frame}

\begin{frame}{Censo}
\textbf{Definición formal:} estudio en que se observa a todos los elementos de la población.
\begin{itemize}
  \item \textbf{Intuición:} no se toma una parte: se mide todo.
  \item \pregunta{¿Medí absolutamente todos los valores de interés?}
  \item \textbf{Ejemplo:} registrar la edad de todos los alumnos matriculados en la UTFSM.
  \item \diferencia{una muestra puede ser pequeña o grande; un censo coincide con toda la población.}
  \item \errorcomun{llamar censo a una encuesta grande pero incompleta.}
\end{itemize}
\end{frame}

\begin{frame}{Relación jerárquica: universo, variable, población y muestra}
\centering
\begin{tikzpicture}[
  node distance=0.8cm,
  box/.style={rectangle, rounded corners, draw=blue!60, fill=blue!8, align=center, minimum width=3.2cm, minimum height=0.75cm},
  smallbox/.style={rectangle, rounded corners, draw=gray!70, fill=gray!8, align=center, minimum width=2.2cm, minimum height=0.55cm},
  arrow/.style={-{Latex}, thick}
]
\node[box] (u) {Universo\\individuos};
\node[box, below left=of u] (v) {Variable\\característica};
\node[box, below right=of u] (p) {Población\\todos los valores};
\node[smallbox, below=of p] (m) {Muestra\\parte medida};
\node[smallbox, below=of m] (c) {Censo\\muestra = población};
\draw[arrow] (u) -- (v);
\draw[arrow] (u) -- (p);
\draw[arrow] (p) -- (m);
\draw[arrow] (m) -- (c);
\end{tikzpicture}
\vspace{0.5em}

\textbf{Clave:} primero se define a quiénes se estudia, luego qué característica se mide y finalmente qué valores se observan.
\end{frame}

\begin{frame}{Dato, información y conocimiento}
\begin{tabular}{>{\bfseries}p{0.22\textwidth}p{0.32\textwidth}p{0.32\textwidth}}
\toprule
Concepto & Intuición & Ejemplo \\
\midrule
Dato & Valor observado individual & 20 años \\
Información & Dato interpretado y con significado & La edad promedio del curso es 20,8 años \\
Conocimiento & Información combinada con experiencia y contexto & El curso es joven; conviene usar ejemplos de primeros años \\
\bottomrule
\end{tabular}
\vspace{0.6em}
\begin{itemize}
  \item \pregunta{¿Estoy viendo un valor, una interpretación o una conclusión contextual?}
  \item \errorcomun{creer que muchos datos equivalen automáticamente a conocimiento.}
\end{itemize}
\end{frame}

\section{Método estadístico}

\begin{frame}{Diseño del estudio}
\textbf{Objetivo:} preparar el estudio antes de recolectar datos.
\begin{itemize}
  \item Definir el problema.
  \item Formular hipótesis o preguntas de investigación.
  \item Determinar universo, variable, población y muestra.
  \item Elegir instrumentos y criterios de medición.
\end{itemize}
\begin{block}{Pregunta guía}
¿Qué quiero responder y qué datos necesito para responderlo?
\end{block}
\errorcomun{empezar recolectando datos sin saber qué decisión o hipótesis se quiere evaluar.}
\end{frame}

\begin{frame}{Ejecución del estudio}
\textbf{Objetivo:} producir y comunicar resultados a partir del diseño.
\begin{enumerate}
  \item Recolección de datos.
  \item Depuración y tabulación.
  \item Construcción de gráficos.
  \item Cálculo de estadísticos.
  \item Análisis e interpretación.
\end{enumerate}
\begin{itemize}
  \item \diferencia{tabular organiza; graficar visualiza; analizar interpreta.}
  \item \errorcomun{presentar tablas o gráficos sin explicar qué significan.}
\end{itemize}
\end{frame}

\section{Clasificación de variables}

\begin{frame}{Variables cualitativas}
\textbf{Definición formal:} variables que describen atributos, categorías o cualidades no numéricas.
\begin{itemize}
  \item \textbf{Intuición:} clasifican a los individuos en grupos.
  \item \pregunta{¿A qué categoría pertenece?}
  \item \textbf{Ejemplos:} sexo, estado civil, color, carrera.
  \item \diferencia{sus valores no representan cantidades medibles, aunque puedan codificarse con números.}
  \item \errorcomun{pensar que una variable es cuantitativa solo porque se codifica como 1, 2 o 3.}
\end{itemize}
\end{frame}

\begin{frame}{Variables cuantitativas}
\textbf{Definición formal:} variables cuyos valores representan cantidades numéricas sobre las que tiene sentido operar aritméticamente.
\begin{itemize}
  \item \textbf{Intuición:} miden cuánto hay de algo.
  \item \pregunta{¿Cuánto mide, pesa, cuesta o cuenta?}
  \item \textbf{Ejemplos:} edad, peso, estatura, ingreso, número de hijos.
  \item \diferencia{a diferencia de las cualitativas, sus valores expresan magnitud.}
  \item \errorcomun{confundir números identificadores, como RUT o código postal, con variables cuantitativas.}
\end{itemize}
\end{frame}

\begin{frame}{Cualitativa vs. cuantitativa}
\begin{tabular}{p{0.24\textwidth}p{0.33\textwidth}p{0.33\textwidth}}
\toprule
Criterio & Cualitativa & Cuantitativa \\
\midrule
Qué expresa & Categoría o atributo & Cantidad o magnitud \\
Pregunta & ¿Cuál tipo? & ¿Cuánto? \\
Ejemplos & carrera, color & edad, ingreso \\
Operaciones & contar categorías & sumar, promediar, comparar magnitudes \\
\bottomrule
\end{tabular}
\vspace{0.6em}
\begin{block}{Regla rápida}
Si el número es una etiqueta, es cualitativa; si mide cantidad, es cuantitativa.
\end{block}
\end{frame}

\begin{frame}{Variable binaria o dicotómica}
\textbf{Definición formal:} variable que solo puede tomar dos valores posibles.
\begin{itemize}
  \item \textbf{Intuición:} responde una alternativa de dos opciones.
  \item \pregunta{¿Sí o no? ¿Presente o ausente?}
  \item \textbf{Ejemplos:} aprobado/reprobado, vivo/muerto, sí/no.
  \item \diferencia{puede ser cualitativa, aunque a veces se codifique como 0 y 1.}
  \item \errorcomun{tratar automáticamente toda variable 0/1 como cuantitativa continua.}
\end{itemize}
\end{frame}

\begin{frame}{Variable discreta}
\textbf{Definición formal:} variable cuantitativa que toma valores contables, usualmente enteros.
\begin{itemize}
  \item \textbf{Intuición:} se cuenta, no se mide con infinitos decimales.
  \item \pregunta{¿Cuántos elementos hay?}
  \item \textbf{Ejemplos:} número de hijos, número de alumnos, cantidad de libros.
  \item \diferencia{entre dos valores consecutivos no hay infinitos valores posibles dentro de la escala.}
  \item \errorcomun{creer que discreta significa ``pocos valores''; lo esencial es que sean contables.}
\end{itemize}
\end{frame}

\begin{frame}{Variable continua}
\textbf{Definición formal:} variable cuantitativa que puede tomar infinitos valores dentro de un intervalo.
\begin{itemize}
  \item \textbf{Intuición:} se mide y puede tener decimales según la precisión del instrumento.
  \item \pregunta{¿Cuánto mide exactamente dentro de un rango?}
  \item \textbf{Ejemplos:} peso, estatura, temperatura, tiempo.
  \item \diferencia{la continua admite valores intermedios; la discreta se basa en conteos.}
  \item \errorcomun{pensar que una medición redondeada deja de ser continua.}
\end{itemize}
\end{frame}

\section{Escalas de medición}

\begin{frame}{Escala nominal}
\textbf{Definición formal:} escala que clasifica datos en categorías sin orden natural.
\begin{itemize}
  \item \textbf{Intuición:} solo pone nombres o etiquetas.
  \item \pregunta{¿A qué grupo pertenece?}
  \item \textbf{Ejemplo:} grupo sanguíneo, carrera, color de ojos.
  \item \diferencia{no permite decir que una categoría es mayor o mejor que otra.}
  \item \errorcomun{ordenar alfabéticamente categorías y creer que ese orden tiene significado estadístico.}
\end{itemize}
\end{frame}

\begin{frame}{Escala ordinal}
\textbf{Definición formal:} escala que clasifica categorías con un orden, pero sin distancias numéricas comparables garantizadas.
\begin{itemize}
  \item \textbf{Intuición:} permite ordenar, no medir diferencias exactas.
  \item \pregunta{¿Cuál está antes, después, más alto o más bajo?}
  \item \textbf{Ejemplo:} malo, regular, bueno, excelente.
  \item \diferencia{ordinal tiene orden; nominal no. Pero ordinal no asegura que los saltos sean iguales.}
  \item \errorcomun{promediar categorías ordinales codificadas como si fueran distancias exactas.}
\end{itemize}
\end{frame}

\begin{frame}{Escala de intervalo}
\textbf{Definición formal:} escala numérica donde las diferencias tienen significado, pero el cero no representa ausencia absoluta.
\begin{itemize}
  \item \textbf{Intuición:} se pueden comparar distancias, no proporciones absolutas.
  \item \pregunta{¿Cuánto difieren dos valores?}
  \item \textbf{Ejemplo:} temperatura en grados Celsius.
  \item \diferencia{20 $^\circ$C no es ``el doble de calor'' que 10 $^\circ$C, porque el cero Celsius no es ausencia de temperatura.}
  \item \errorcomun{interpretar razones cuando solo son válidas diferencias.}
\end{itemize}
\end{frame}

\begin{frame}{Escala de razón}
\textbf{Definición formal:} escala numérica con diferencias significativas y cero absoluto.
\begin{itemize}
  \item \textbf{Intuición:} permite decir ``el doble'', ``la mitad'' o ``ausencia de''.
  \item \pregunta{¿Cuánto y en qué proporción?}
  \item \textbf{Ejemplos:} masa, distancia, tiempo, ingreso.
  \item \diferencia{a diferencia del intervalo, el cero sí representa ausencia de la magnitud.}
  \item \errorcomun{confundir intervalo y razón solo porque ambas son numéricas.}
\end{itemize}
\end{frame}

\begin{frame}{Resumen de escalas de medición}
\centering
\begin{tabular}{p{0.19\textwidth}p{0.18\textwidth}p{0.18\textwidth}p{0.24\textwidth}}
\toprule
Escala & Clasifica & Ordena & Permite \\
\midrule
Nominal & Sí & No & contar categorías \\
Ordinal & Sí & Sí & comparar posiciones \\
Intervalo & Sí & Sí & comparar diferencias \\
Razón & Sí & Sí & comparar diferencias y proporciones \\
\bottomrule
\end{tabular}
\vspace{0.6em}
\begin{block}{Jerarquía}
Razón $\supset$ intervalo $\supset$ ordinal $\supset$ nominal, en cantidad de operaciones interpretables.
\end{block}
\end{frame}

\section{Dimensionalidad}

\begin{frame}{Dimensionalidad de los datos}
\begin{itemize}
  \item \textbf{Unidimensional:} se estudia una sola variable. Ejemplo: edad.
  \item \textbf{Bidimensional:} se estudian dos variables simultáneamente. Ejemplo: edad y peso.
  \item \textbf{Multidimensional:} se estudian tres o más variables. Ejemplo: edad, peso, ingreso y carrera.
\end{itemize}
\begin{block}{Pregunta guía}
¿Cuántas variables se analizan al mismo tiempo?
\end{block}
\errorcomun{confundir tamaño de muestra con dimensionalidad: muchos datos no implican muchas variables.}
\end{frame}

\section{Estadística descriptiva}

\begin{frame}{Objetivos de la estadística descriptiva}
\begin{columns}[T]
\column{0.5\textwidth}
\begin{itemize}
  \item Organizar datos.
  \item Resumir datos.
  \item Presentar datos.
  \item Visualizar datos.
\end{itemize}
\column{0.45\textwidth}
\begin{block}{Herramientas}
Tablas, gráficos y estadísticos como media, mediana, moda, rango y desviación estándar.
\end{block}
\end{columns}
\vspace{0.5em}
\pregunta{¿Cómo puedo entender rápidamente el comportamiento de los datos observados?}
\end{frame}

\section{Tablas de frecuencia}

\begin{frame}{Tabla de frecuencia: idea general}
\textbf{Definición formal:} tabla que resume cuántas veces aparece cada valor o categoría y, si corresponde, sus proporciones y acumulados.
\begin{itemize}
  \item \textbf{Intuición:} transforma una lista larga de datos en un resumen ordenado.
  \item \pregunta{¿Cómo se distribuyen los datos entre valores o categorías?}
  \item \textbf{Ejemplo:} contar cuántos estudiantes obtienen cada nota o pertenecen a cada carrera.
  \item \diferencia{una tabla de frecuencia resume una variable; una tabla de contingencia cruza dos o más variables.}
\end{itemize}
\end{frame}

\begin{frame}{Frecuencia absoluta: $f_i$}
\textbf{Definición formal:} número de veces que aparece el valor o categoría $i$.
\begin{itemize}
  \item \textbf{Intuición:} conteo directo.
  \item \pregunta{¿Cuántas veces apareció?}
  \item \textbf{Ejemplo:} si 8 estudiantes tienen nota 5, entonces $f_i=8$ para la nota 5.
  \item \diferencia{es conteo, no proporción.}
  \item \errorcomun{comparar frecuencias absolutas entre grupos de distinto tamaño sin considerar proporciones.}
\end{itemize}
\end{frame}

\begin{frame}{Frecuencia relativa: $h_i$}
\textbf{Definición formal:} proporción que representa la frecuencia absoluta respecto del total: $h_i = f_i/n$.
\begin{itemize}
  \item \textbf{Intuición:} parte del total que corresponde a un valor o categoría.
  \item \pregunta{¿Qué fracción o porcentaje representa?}
  \item \textbf{Ejemplo:} si $f_i=8$ y $n=40$, entonces $h_i=8/40=0{,}20=20\%$.
  \item \diferencia{la absoluta cuenta; la relativa divide por el total.}
  \item \errorcomun{olvidar que la suma de frecuencias relativas debe ser 1, o 100\%.}
\end{itemize}
\end{frame}

\begin{frame}{Frecuencia absoluta acumulada: $F_i$}
\textbf{Definición formal:} suma de frecuencias absolutas hasta el valor o clase $i$: $F_i=f_1+\cdots+f_i$.
\begin{itemize}
  \item \textbf{Intuición:} conteo acumulado desde el inicio de la tabla.
  \item \pregunta{¿Cuántos datos van acumulados hasta aquí?}
  \item \textbf{Ejemplo:} si hasta nota 4 se acumulan 25 estudiantes, entonces $F_i=25$.
  \item \diferencia{acumulada considera valores anteriores; $f_i$ solo mira la fila actual.}
  \item \errorcomun{acumular categorías nominales sin un orden significativo.}
\end{itemize}
\end{frame}

\begin{frame}{Frecuencia relativa acumulada: $H_i$}
\textbf{Definición formal:} suma de frecuencias relativas hasta el valor o clase $i$: $H_i=h_1+\cdots+h_i$.
\begin{itemize}
  \item \textbf{Intuición:} porcentaje acumulado.
  \item \pregunta{¿Qué porcentaje va acumulado hasta aquí?}
  \item \textbf{Ejemplo:} si $H_i=0{,}75$, entonces el 75\% de los datos está en ese valor o por debajo de él.
  \item \diferencia{$F_i$ acumula conteos; $H_i$ acumula proporciones.}
  \item \errorcomun{interpretar $H_i$ sin revisar el orden de la variable.}
\end{itemize}
\end{frame}

\begin{frame}{Ejemplo de tabla de frecuencia}
\centering
\begin{tabular}{ccccc}
\toprule
Nota & $f_i$ & $h_i$ & $F_i$ & $H_i$ \\
\midrule
3 & 4 & 0,10 & 4 & 0,10 \\
4 & 10 & 0,25 & 14 & 0,35 \\
5 & 16 & 0,40 & 30 & 0,75 \\
6 & 8 & 0,20 & 38 & 0,95 \\
7 & 2 & 0,05 & 40 & 1,00 \\
\bottomrule
\end{tabular}
\vspace{0.6em}
\begin{itemize}
  \item Total de datos: $n=40$.
  \item La última frecuencia absoluta acumulada debe ser $n$.
  \item La última frecuencia relativa acumulada debe ser 1 o 100\%.
\end{itemize}
\end{frame}

\section{Variables continuas e intervalos}

\begin{frame}{Rango o recorrido}
\textbf{Definición formal:} diferencia entre el valor máximo y el valor mínimo: $R=x_{\max}-x_{\min}$.
\begin{itemize}
  \item \textbf{Intuición:} extensión total que cubren los datos.
  \item \pregunta{¿Qué tan separados están los extremos?}
  \item \textbf{Ejemplo:} si la estatura mínima es 150 cm y la máxima 190 cm, entonces $R=40$ cm.
  \item \diferencia{el rango describe extremos; no indica cómo se distribuyen los datos internamente.}
  \item \errorcomun{concluir que dos conjuntos son similares solo porque tienen el mismo rango.}
\end{itemize}
\end{frame}

\begin{frame}{Intervalo de clase}
\textbf{Definición formal:} grupo o tramo de valores usado para resumir variables continuas o datos numerosos.
\begin{itemize}
  \item \textbf{Intuición:} una ``caja'' donde caen valores cercanos.
  \item \pregunta{¿En qué tramo cae este dato?}
  \item \textbf{Ejemplo:} estaturas entre 160 y 170 cm.
  \item \diferencia{el intervalo agrupa valores; el dato individual es un valor específico.}
  \item \errorcomun{crear intervalos superpuestos o dejar huecos entre intervalos.}
\end{itemize}
\end{frame}

\begin{frame}{Amplitud del intervalo}
\textbf{Definición formal:} ancho de cada intervalo de clase, usualmente calculado como $A=R/k$, donde $k$ es el número de intervalos.
\begin{itemize}
  \item \textbf{Intuición:} tamaño de cada tramo.
  \item \pregunta{¿Cuánto mide cada intervalo?}
  \item \textbf{Ejemplo:} si $R=40$ y se usan $k=5$ intervalos, entonces $A=8$.
  \item \diferencia{rango mide la extensión total; amplitud mide el ancho de cada intervalo.}
  \item \errorcomun{confundir amplitud con número de intervalos.}
\end{itemize}
\end{frame}

\begin{frame}{Marca de clase}
\textbf{Definición formal:} punto medio de un intervalo: $m_i=(L_i+U_i)/2$, donde $L_i$ es el límite inferior y $U_i$ el superior.
\begin{itemize}
  \item \textbf{Intuición:} valor representativo del intervalo.
  \item \pregunta{¿Qué valor resume este tramo?}
  \item \textbf{Ejemplo:} para el intervalo [160, 170), la marca de clase es 165.
  \item \diferencia{la marca no es un dato observado necesariamente; es un representante del intervalo.}
  \item \errorcomun{creer que todos los datos del intervalo valen exactamente la marca de clase.}
\end{itemize}
\end{frame}

\begin{frame}{Número de intervalos: $\sqrt{n}$ y Sturges}
\textbf{Objetivo:} elegir una cantidad razonable de intervalos para agrupar datos.
\begin{columns}[T]
\column{0.48\textwidth}
\begin{block}{Regla $\sqrt{n}$}
$k \approx \sqrt{n}$
\end{block}
Simple y rápida para una primera aproximación.
\column{0.48\textwidth}
\begin{block}{Regla de Sturges}
$k \approx 1+\log_2(n)$
\end{block}
Muy usada para construir histogramas y tablas agrupadas.
\end{columns}
\vspace{0.6em}
\errorcomun{tratar estas reglas como verdades absolutas; son criterios prácticos, no leyes.}
\end{frame}

\begin{frame}{Ejemplo: construir intervalos}
\begin{itemize}
  \item Datos de estatura: mínimo 150 cm, máximo 190 cm, $n=25$.
  \item Rango: $R=190-150=40$ cm.
  \item Número de intervalos por $\sqrt{n}$: $k\approx\sqrt{25}=5$.
  \item Amplitud: $A=R/k=40/5=8$ cm.
  \item Primeros intervalos: [150,158), [158,166), [166,174), \dots
\end{itemize}
\begin{block}{Chequeo}
Los intervalos deben cubrir todo el rango sin superponerse ni dejar valores fuera.
\end{block}
\end{frame}

\section{Tablas de contingencia}

\begin{frame}{Tabla de contingencia}
\textbf{Definición formal:} tabla que cruza dos o más variables categóricas para estudiar la distribución conjunta y posible asociación entre ellas.
\begin{itemize}
  \item \textbf{Intuición:} muestra cómo se combinan categorías de dos variables.
  \item \pregunta{¿Cómo se relacionan las categorías de una variable con las de otra?}
  \item \textbf{Ejemplo:} carrera del estudiante versus aprobación del curso.
  \item \diferencia{tabla de frecuencia: una variable; contingencia: dos o más variables.}
  \item \errorcomun{concluir causalidad solo porque aparece asociación en la tabla.}
\end{itemize}
\end{frame}

\begin{frame}{Elementos de una tabla de contingencia}
\begin{itemize}
  \item \textbf{Variable fila:} variable ubicada en las filas de la tabla.
  \item \textbf{Variable columna:} variable ubicada en las columnas.
  \item \textbf{Categorías:} valores posibles de cada variable.
  \item \textbf{Frecuencias conjuntas:} conteos dentro de cada celda.
  \item \textbf{Totales marginales:} sumas por fila o columna.
\end{itemize}
\begin{block}{Pregunta clave}
¿La distribución de una variable cambia según la categoría de la otra?
\end{block}
\end{frame}

\begin{frame}{Ejemplo de tabla de contingencia}
\centering
\begin{tabular}{lccc}
\toprule
Carrera & Aprobó & Reprobó & Total \\
\midrule
Ingeniería & 45 & 10 & 55 \\
Arquitectura & 20 & 8 & 28 \\
Comercial & 30 & 12 & 42 \\
\midrule
Total & 95 & 30 & 125 \\
\bottomrule
\end{tabular}
\vspace{0.7em}
\begin{itemize}
  \item Variable fila: carrera.
  \item Variable columna: resultado del curso.
  \item Las celdas muestran frecuencias conjuntas.
  \item Para comparar carreras conviene usar porcentajes por fila, no solo conteos absolutos.
\end{itemize}
\end{frame}

\section{Cierre}

\begin{frame}{Errores comunes que debes evitar}
\begin{itemize}
  \item Confundir individuos del universo con valores de la población.
  \item Llamar cuantitativa a una variable solo porque está codificada con números.
  \item Usar acumulados en variables nominales sin orden natural.
  \item Confundir rango, amplitud, intervalo y marca de clase.
  \item Comparar frecuencias absolutas entre grupos de distinto tamaño.
  \item Interpretar asociación en tablas de contingencia como causalidad automática.
\end{itemize}
\end{frame}

\begin{frame}{Checklist de estudio rápido}
\begin{enumerate}
  \item ¿Sé distinguir universo, variable, población, muestra y censo?
  \item ¿Puedo explicar la diferencia entre dato, información y conocimiento?
  \item ¿Clasifico variables por naturaleza y por cantidad de valores posibles?
  \item ¿Reconozco las escalas nominal, ordinal, intervalo y razón?
  \item ¿Calculo e interpreto $f_i$, $h_i$, $F_i$ y $H_i$?
  \item ¿Construyo intervalos usando rango, amplitud y marca de clase?
  \item ¿Diferencio tabla de frecuencia y tabla de contingencia?
\end{enumerate}
\end{frame}

\begin{frame}{Idea final}
\centering
\Large
La estadística descriptiva no solo resume datos: ayuda a transformar observaciones en información clara para tomar mejores decisiones.
\end{frame}


\section{Unidad 2: Estadística Descriptiva Avanzada}

\begin{frame}{Unidad 2: mapa de contenidos}
\begin{columns}[T]
\column{0.48\textwidth}
\begin{enumerate}
  \item Representación gráfica de datos
  \item Medidas de resumen
  \item Medidas de posición
\end{enumerate}
\column{0.48\textwidth}
\begin{enumerate}\setcounter{enumi}{3}
  \item Percentiles y cuantiles
  \item Medidas de dispersión
  \item Resúmenes visuales por bloque
\end{enumerate}
\end{columns}
\vspace{0.5em}
\begin{block}{Continuidad con la unidad anterior}
Las tablas ordenan los datos; los gráficos y medidas de resumen ayudan a interpretar forma, centro, posición y variabilidad.
\end{block}
\end{frame}

\subsection{Representación gráfica de datos}

\begin{frame}{Representación gráfica de datos}
\textbf{Definición formal:} recurso visual que representa datos mediante posiciones, longitudes, áreas, sectores, símbolos o puntos.
\begin{itemize}
  \item \textbf{Intuición:} permite ver patrones que una tabla puede ocultar.
  \item \pregunta{¿Qué forma, concentración, tendencia o relación muestran los datos?}
  \item \diferencia{la tabla entrega valores exactos; el gráfico facilita comparación visual e interpretación global.}
  \item \errorcomun{usar gráficos decorativos que distorsionan proporciones o dificultan comparar magnitudes.}
\end{itemize}
\end{frame}

\begin{frame}{Cuándo usar gráficos en lugar de tablas}
\begin{itemize}
  \item Usa \textbf{tablas} cuando necesites valores exactos, acumulados o cálculos posteriores.
  \item Usa \textbf{gráficos} cuando quieras comparar categorías, detectar forma, observar acumulados o comunicar resultados rápidamente.
  \item En reportes claros, tabla y gráfico se complementan: una muestra precisión y el otro muestra patrón.
\end{itemize}
\begin{block}{Limitación clave}
Un gráfico puede sugerir una interpretación, pero no reemplaza revisar escala, total de datos, unidades y método de construcción.
\end{block}
\end{frame}

\begin{frame}{Gráfico de barras}
\textbf{Definición formal:} gráfico que representa frecuencias de valores discretos o categorías mediante barras separadas.
\begin{itemize}
  \item \textbf{Intuición:} cada barra es un conteo o porcentaje de una categoría.
  \item \pregunta{¿Qué categoría o valor discreto aparece más o menos?}
  \item \diferencia{en barras hay separación entre categorías; en histogramas las barras son contiguas porque representan intervalos continuos.}
  \item \errorcomun{ordenar barras de una variable ordinal como si fueran nominales, perdiendo el sentido del orden.}
\end{itemize}
\end{frame}

\begin{frame}{Ejemplo gráfico: barras y polígono de frecuencias}
\centering
\begin{tikzpicture}[x=1.25cm,y=0.18cm]
  \draw[->] (0,0) -- (6.2,0) node[right]{Nota};
  \draw[->] (0,0) -- (0,18) node[above]{$f_i$};
  \foreach \x/\lab in {1/3,2/4,3/5,4/6,5/7}{\node[below] at (\x,0) {\lab};}
  \foreach \y in {5,10,15}{\draw (-0.07,\y) -- (0.07,\y) node[left=3pt]{\y};}
  \foreach \x/\h in {1/4,2/10,3/16,4/8,5/2}{\fill[blue!35] (\x-0.3,0) rectangle (\x+0.3,\h);}
  \draw[red,thick] (1,4) -- (2,10) -- (3,16) -- (4,8) -- (5,2);
  \foreach \x/\h in {1/4,2/10,3/16,4/8,5/2}{\fill[red] (\x,\h) circle (2pt);}
\end{tikzpicture}
\vspace{0.4em}
\begin{itemize}
  \item Ejemplo completo: de $n=40$, la nota 5 tiene $f_i=16$ y $h_i=16/40=40\%$.
  \item La línea roja une puntos de frecuencia y funciona como polígono de frecuencias.
\end{itemize}
\end{frame}

\begin{frame}{Gráfico sectorial o de torta}
\textbf{Definición formal:} gráfico circular dividido en sectores proporcionales a frecuencias relativas o porcentajes.
\begin{itemize}
  \item \textbf{Intuición:} muestra cómo se reparte un total entre categorías.
  \item \pregunta{¿Qué parte del total representa cada categoría?}
  \item \diferencia{sirve para composición porcentual; no es ideal para comparar muchas categorías parecidas.}
  \item \errorcomun{usar tortas con demasiadas categorías o sin indicar porcentajes, haciendo difícil comparar sectores.}
\end{itemize}
\end{frame}

\begin{frame}{Ejemplo gráfico: sectores proporcionales}
\centering
\begin{tikzpicture}[scale=1.05]
  \fill[blue!35] (0,0) -- (0:1.8) arc (0:144:1.8) -- cycle;
  \fill[green!35] (0,0) -- (144:1.8) arc (144:234:1.8) -- cycle;
  \fill[orange!45] (0,0) -- (234:1.8) arc (234:306:1.8) -- cycle;
  \fill[red!35] (0,0) -- (306:1.8) arc (306:360:1.8) -- cycle;
  \draw (0,0) circle (1.8);
  \node at (70:1.05) {40\%};
  \node at (188:1.05) {25\%};
  \node at (270:1.05) {20\%};
  \node at (335:1.05) {15\%};
  \node[right] at (2.4,0.9) {\textcolor{blue!70}{\rule{0.6cm}{0.18cm}} Nota 5};
  \node[right] at (2.4,0.45) {\textcolor{green!70}{\rule{0.6cm}{0.18cm}} Nota 4};
  \node[right] at (2.4,0.0) {\textcolor{orange!80}{\rule{0.6cm}{0.18cm}} Nota 6};
  \node[right] at (2.4,-0.45) {\textcolor{red!70}{\rule{0.6cm}{0.18cm}} Otras};
\end{tikzpicture}
\vspace{0.4em}
\begin{itemize}
  \item Ejemplo: si 16 de 40 estudiantes obtienen nota 5, el sector es $40\%$ del círculo.
  \item Ángulo del sector: $0{,}40\times 360^\circ=144^\circ$.
\end{itemize}
\end{frame}

\begin{frame}{Histograma}
\textbf{Definición formal:} gráfico para variables continuas agrupadas en intervalos, donde barras contiguas representan frecuencias por clase.
\begin{itemize}
  \item \textbf{Intuición:} muestra la forma aproximada de la distribución.
  \item \pregunta{¿En qué intervalos se concentran los valores?}
  \item \diferencia{usa intervalos y barras pegadas; el gráfico de barras usa categorías o valores discretos separados.}
  \item \errorcomun{tratar la marca de clase como si fuera el valor real de todos los datos del intervalo.}
\end{itemize}
\end{frame}

\begin{frame}{Ejemplo gráfico: histograma}
\centering
\begin{tikzpicture}[x=0.08cm,y=0.18cm]
  \draw[->] (148,0) -- (194,0) node[right]{Estatura};
  \draw[->] (150,0) -- (150,17) node[above]{$f_i$};
  \foreach \x in {150,158,166,174,182,190}{\draw (\x,0) -- (\x,-0.4) node[below]{\x};}
  \foreach \y in {5,10,15}{\draw (149.2,\y)--(150.8,\y) node[left=5pt]{\y};}
  \fill[purple!35] (150,0) rectangle (158,3);
  \fill[purple!35] (158,0) rectangle (166,7);
  \fill[purple!35] (166,0) rectangle (174,12);
  \fill[purple!35] (174,0) rectangle (182,8);
  \fill[purple!35] (182,0) rectangle (190,5);
  \draw (150,0) rectangle (158,3);
  \draw (158,0) rectangle (166,7);
  \draw (166,0) rectangle (174,12);
  \draw (174,0) rectangle (182,8);
  \draw (182,0) rectangle (190,5);
\end{tikzpicture}
\vspace{0.4em}
\begin{itemize}
  \item Intervalo con mayor frecuencia: [166,174), con $f_i=12$.
  \item Marca de clase de [166,174): $(166+174)/2=170$.
\end{itemize}
\end{frame}

\begin{frame}{Pictograma}
\textbf{Definición formal:} representación gráfica que usa símbolos o imágenes para indicar frecuencias según una escala definida.
\begin{itemize}
  \item \textbf{Intuición:} reemplaza barras por íconos contables.
  \item \pregunta{¿Cuántos casos representa visualmente cada símbolo?}
  \item \diferencia{un pictograma informa mediante una escala; un gráfico decorado solo agrega imágenes sin valor estadístico.}
  \item \errorcomun{cambiar el tamaño del símbolo en alto y ancho, exagerando visualmente el área representada.}
\end{itemize}
\end{frame}

\begin{frame}{Ejemplo gráfico: pictograma}
\centering
\begin{tabular}{ll}
\toprule
Categoría & Pictograma \\
\midrule
Sección A & $\blacksquare\ \blacksquare\ \blacksquare\ \blacksquare$ \\
Sección B & $\blacksquare\ \blacksquare\ \blacksquare$ \\
Sección C & $\blacksquare\ \blacksquare$ \\
\bottomrule
\end{tabular}
\vspace{0.7em}
\begin{block}{Escala}
Cada $\blacksquare$ representa 5 estudiantes.
\end{block}
\begin{itemize}
  \item Sección A: $4\times 5=20$ estudiantes.
  \item Sección C: $2\times 5=10$ estudiantes.
\end{itemize}
\end{frame}

\begin{frame}{Ojiva}
\textbf{Definición formal:} gráfico de frecuencias acumuladas, absolutas o relativas, usualmente construido sobre límites superiores de intervalos.
\begin{itemize}
  \item \textbf{Intuición:} muestra cuánto se ha acumulado hasta cierto umbral.
  \item \pregunta{¿Qué porcentaje o cuántos datos están por debajo de un valor?}
  \item \diferencia{el histograma muestra frecuencias por intervalo; la ojiva muestra acumulados.}
  \item \errorcomun{interpretar un punto de la ojiva como frecuencia simple y no acumulada.}
\end{itemize}
\end{frame}

\begin{frame}{Ejemplo gráfico: ojiva}
\centering
\begin{tikzpicture}[x=0.08cm,y=0.04cm]
  \draw[->] (148,0) -- (194,0) node[right]{Límite superior};
  \draw[->] (150,0) -- (150,105) node[above]{$H_i$};
  \foreach \x in {150,158,166,174,182,190}{\draw (\x,0)--(\x,-2) node[below]{\x};}
  \foreach \y in {25,50,75,100}{\draw (149.2,\y)--(150.8,\y) node[left=5pt]{\y\%};}
  \draw[blue,thick] (150,0) -- (158,9) -- (166,29) -- (174,63) -- (182,86) -- (190,100);
  \foreach \x/\y in {150/0,158/9,166/29,174/63,182/86,190/100}{\fill[blue] (\x,\y) circle (2pt);}
  \draw[dashed] (174,0)--(174,63)--(150,63);
\end{tikzpicture}
\vspace{0.4em}
\begin{itemize}
  \item Interpretación: cerca del 63\% de las estaturas está bajo 174 cm.
  \item Sirve para ubicar percentiles mediante interpolación visual o numérica.
\end{itemize}
\end{frame}

\begin{frame}{Diagrama de tallo y hoja}
\textbf{Definición formal:} representación que separa cada dato en tallo y hoja para ordenar observaciones conservando los datos originales.
\begin{itemize}
  \item \textbf{Intuición:} combina tabla ordenada e idea visual de distribución.
  \item \pregunta{¿Cómo se distribuyen los datos sin perder los valores exactos?}
  \item \diferencia{a diferencia del histograma, permite reconstruir los datos originales.}
  \item \errorcomun{no indicar la clave de lectura, dejando ambiguo si $7\mid 2$ significa 72, 7,2 u otro valor.}
\end{itemize}
\end{frame}

\begin{frame}{Ejemplo: tallo y hoja}
\centering
\begin{tabular}{c|l}
\toprule
Tallo & Hojas \\
\midrule
5 & 8\quad 9 \\
6 & 0\quad 2\quad 4\quad 7 \\
7 & 1\quad 2\quad 5\quad 8 \\
8 & 0\quad 3 \\
\bottomrule
\end{tabular}
\vspace{0.7em}
\begin{itemize}
  \item Clave: $7\mid 2$ representa 72.
  \item Datos reconstruidos: 58, 59, 60, 62, 64, 67, 71, 72, 75, 78, 80, 83.
  \item Se observa concentración entre 60 y 78.
\end{itemize}
\end{frame}

\begin{frame}{Diagrama de dispersión}
\textbf{Definición formal:} gráfico de pares ordenados $(x,y)$ para estudiar visualmente la relación entre dos variables cuantitativas.
\begin{itemize}
  \item \textbf{Intuición:} cada punto representa un individuo medido en dos variables.
  \item \pregunta{¿Existe asociación positiva, negativa, nula o valores atípicos?}
  \item \diferencia{muestra relación entre dos variables; no resume una sola distribución como un histograma.}
  \item \errorcomun{confundir correlación visual con causalidad demostrada.}
\end{itemize}
\end{frame}

\begin{frame}{Ejemplo gráfico: diagrama de dispersión}
\centering
\begin{tikzpicture}[x=0.55cm,y=0.45cm]
  \draw[->] (0,0) -- (9.5,0) node[right]{Horas de estudio};
  \draw[->] (0,0) -- (0,9.5) node[above]{Nota};
  \foreach \x in {2,4,6,8}{\draw (\x,0)--(\x,-0.1) node[below]{\x};}
  \foreach \y in {2,4,6,8}{\draw (0,\y)--(-0.1,\y) node[left]{\y};}
  \foreach \x/\y in {1/2,2/3,3/3.8,4/5,5/5.4,6/6.5,7/7.2,8/8.1,2/7.8}{\fill[blue!70] (\x,\y) circle (2.5pt);}
  \draw[red,dashed] (1,2.2)--(8,8);
  \node[red,right] at (2,8.1) {posible atípico};
\end{tikzpicture}
\vspace{0.4em}
\begin{itemize}
  \item La nube sugiere asociación positiva: más horas, mayor nota.
  \item El punto alto con pocas horas podría ser un valor atípico o un caso especial.
\end{itemize}
\end{frame}

\begin{frame}{Resumen visual: gráficos}
\small
\begin{tabular}{p{0.22\textwidth}p{0.33\textwidth}p{0.33\textwidth}}
\toprule
Gráfico & Úsalo cuando & Evita confundirlo con \\
\midrule
Barras & categorías o discretas & histograma \\
Torta & composición porcentual & comparación precisa entre muchas categorías \\
Histograma & variable continua agrupada & barras separadas \\
Ojiva & acumulados y percentiles & frecuencia simple \\
Tallo y hoja & pocos datos y valores exactos & tabla de intervalos \\
Dispersión & relación entre dos variables & prueba de causalidad \\
\bottomrule
\end{tabular}
\end{frame}

\subsection{Medidas de resumen y posición}

\begin{frame}{Medidas de resumen}
\textbf{Definición formal:} valores numéricos que sintetizan una distribución en aspectos como centro, posición, dispersión o forma.
\begin{itemize}
  \item \textbf{Intuición:} reducen muchos datos a pocos indicadores interpretables.
  \item \pregunta{¿Cómo describo la distribución sin listar todos los datos?}
  \item \diferencia{un gráfico muestra forma visual; una medida entrega un número preciso.}
  \item \errorcomun{usar una sola medida, como la media, para describir todo un conjunto de datos.}
\end{itemize}
\end{frame}

\begin{frame}{Parámetros y estadísticos}
\begin{columns}[T]
\column{0.48\textwidth}
\textbf{Parámetro}
\begin{itemize}
  \item Resume una población.
  \item Suele representarse con letras griegas.
  \item Ejemplo: media poblacional $\mu$.
\end{itemize}
\column{0.48\textwidth}
\textbf{Estadístico}
\begin{itemize}
  \item Resume una muestra.
  \item Se calcula con datos observados.
  \item Ejemplo: media muestral $\bar{x}$.
\end{itemize}
\end{columns}
\vspace{0.4em}
\diferencia{el parámetro describe el total poblacional; el estadístico describe la muestra y puede estimar al parámetro.}
\end{frame}

\begin{frame}{Concepto de posición estadística}
\textbf{Definición formal:} ubicación de un valor dentro de la distribución, ya sea en el centro o en una posición relativa específica.
\begin{itemize}
  \item \textbf{Tendencia central:} media, mediana y moda buscan representar el centro o valor típico.
  \item \textbf{Posición no central:} cuartiles, deciles y percentiles ubican cortes de la distribución.
  \item \pregunta{¿Dónde se ubica un dato o un resumen dentro del conjunto?}
  \item \errorcomun{creer que todas las medidas de posición describen solo el centro.}
\end{itemize}
\end{frame}

\begin{frame}{Media aritmética: datos no agrupados}
\textbf{Definición formal:} suma de todos los datos dividida por la cantidad de observaciones.
\[
\bar{x}=\frac{x_1+x_2+\cdots+x_n}{n}
\]
\begin{itemize}
  \item \textbf{Intuición:} reparto equitativo del total entre todos los datos.
  \item \pregunta{¿Cuál sería el valor si todos tuvieran la misma cantidad?}
  \item \diferencia{usa todos los valores; por eso es sensible a valores atípicos.}
  \item \textbf{Ejemplo:} datos 4, 5, 6, 7: $\bar{x}=22/4=5{,}5$.
\end{itemize}
\end{frame}

\begin{frame}{Mediana: datos no agrupados}
\textbf{Definición formal:} valor central de los datos ordenados; divide la distribución en dos mitades.
\begin{itemize}
  \item \textbf{Caso impar:} 3, 4, \textbf{6}, 8, 9 $\Rightarrow Me=6$.
  \item \textbf{Caso par:} 3, 4, 6, 8 $\Rightarrow Me=(4+6)/2=5$.
  \item \pregunta{¿Qué valor deja 50\% de datos a cada lado?}
  \item \diferencia{la mediana depende del orden; la media depende de magnitudes y es más sensible a extremos.}
  \item \errorcomun{calcular la mediana sin ordenar primero los datos.}
\end{itemize}
\end{frame}

\begin{frame}{Moda: datos no agrupados}
\textbf{Definición formal:} valor o categoría que aparece con mayor frecuencia.
\begin{itemize}
  \item \textbf{Intuición:} lo más común.
  \item \pregunta{¿Qué valor se repite más?}
  \item \textbf{Ejemplo:} 2, 3, 3, 4, 5, 5 tiene dos modas: 3 y 5.
  \item \diferencia{puede aplicarse a variables cualitativas; media y mediana requieren más estructura numérica u ordinal.}
  \item \errorcomun{forzar una única moda cuando la distribución puede ser bimodal o multimodal.}
\end{itemize}
\end{frame}

\begin{frame}{Media aritmética ponderada}
\textbf{Definición formal:} promedio donde cada valor tiene un peso asociado.
\[
\bar{x}_p=\frac{\sum w_i x_i}{\sum w_i}
\]
\begin{itemize}
  \item \textbf{Intuición:} algunos valores cuentan más que otros.
  \item \pregunta{¿Cuál es el promedio considerando importancia relativa?}
  \item \textbf{Ejemplo créditos:} notas 6,0; 5,0; 4,0 con créditos 4, 3, 2: $\bar{x}_p=(24+15+8)/9=5{,}22$.
  \item \errorcomun{promediar notas con créditos distintos como si todas pesaran igual.}
\end{itemize}
\end{frame}

\begin{frame}{Media para datos agrupados}
\textbf{Definición formal:} promedio aproximado usando frecuencias y, si hay intervalos, marcas de clase.
\[
\bar{x}=\frac{\sum f_i x_i}{n}\quad\text{o}\quad \bar{x}\approx\frac{\sum f_i m_i}{n}
\]
\begin{itemize}
  \item \textbf{Ejemplo discreto:} notas 4, 5, 6 con frecuencias 2, 5, 3: $\bar{x}=(8+25+18)/10=5{,}1$.
  \item \textbf{Ejemplo intervalos:} usar $m_i$ como representante de cada clase.
  \item \diferencia{en intervalos la media es aproximada porque no se conocen todos los datos originales.}
\end{itemize}
\end{frame}

\begin{frame}{Mediana para datos agrupados}
\textbf{Definición formal:} valor aproximado que ubica el 50\% acumulado dentro del intervalo mediano.
\[
Me=L+\left(\frac{\frac{n}{2}-F_{ant}}{f_m}\right)A
\]
\begin{itemize}
  \item $L$: límite inferior del intervalo mediano; $F_{ant}$: acumulada anterior.
  \item $f_m$: frecuencia del intervalo mediano; $A$: amplitud.
  \item \textbf{Ejemplo:} $n=40$, $L=50$, $F_{ant}=14$, $f_m=16$, $A=10$.
  \item $Me=50+((20-14)/16)10=53{,}75$.
\end{itemize}
\end{frame}

\begin{frame}{Moda para datos agrupados}
\textbf{Definición formal:} aproximación de la moda dentro del intervalo con mayor frecuencia.
\[
Mo=L+\left(\frac{d_1}{d_1+d_2}\right)A
\]
\begin{itemize}
  \item $d_1=f_m-f_{ant}$ y $d_2=f_m-f_{sig}$.
  \item \textbf{Ejemplo:} $L=50$, $A=10$, $f_m=16$, $f_{ant}=10$, $f_{sig}=8$.
  \item $d_1=6$, $d_2=8$, entonces $Mo=50+(6/14)10=54{,}29$.
  \item \errorcomun{aplicar la fórmula si no hay un intervalo modal claro.}
\end{itemize}
\end{frame}

\begin{frame}{Relación entre media, mediana y moda}
\begin{columns}[T]
\column{0.5\textwidth}
\begin{itemize}
  \item \textbf{Simétrica:} media $\approx$ mediana $\approx$ moda.
  \item \textbf{Asimetría positiva:} media $>$ mediana $>$ moda.
  \item \textbf{Asimetría negativa:} media $<$ mediana $<$ moda.
\end{itemize}
\column{0.45\textwidth}
\centering
\begin{tikzpicture}[scale=0.75]
  \draw[->] (0,0)--(5,0);
  \draw[thick,blue] plot[smooth] coordinates {(0.3,0.1) (1.2,1.9) (2.1,0.8) (4.6,0.2)};
  \draw[dashed] (1.2,0)--(1.2,1.8) node[above]{Mo};
  \draw[dashed] (1.8,0)--(1.8,1.0) node[above]{Me};
  \draw[dashed] (2.5,0)--(2.5,0.55) node[above]{$\bar{x}$};
\end{tikzpicture}
\end{columns}
\vspace{0.4em}
\errorcomun{interpretar la media como centro representativo cuando existe fuerte asimetría o valores atípicos.}
\end{frame}

\begin{frame}{Resumen visual: medidas de posición}
\small
\begin{tabular}{p{0.2\textwidth}p{0.36\textwidth}p{0.32\textwidth}}
\toprule
Medida & Responde & Cuidado principal \\
\midrule
Media & ¿valor promedio de equilibrio? & sensible a extremos \\
Mediana & ¿centro por posición? & ordenar antes de calcular \\
Moda & ¿valor más frecuente? & puede no existir o haber varias \\
Media ponderada & ¿promedio con pesos? & pesos deben aplicarse correctamente \\
Cuantiles & ¿ubicación relativa? & convención de cálculo usada \\
\bottomrule
\end{tabular}
\end{frame}

\subsection{Percentiles y cuantiles}

\begin{frame}{Percentiles}
\textbf{Definición formal:} valores que dividen la distribución en 100 partes porcentuales; el percentil $q$ deja aproximadamente $q\%$ de datos por debajo.
\begin{itemize}
  \item \textbf{Intuición:} posición relativa expresada como porcentaje.
  \item \pregunta{¿Qué valor marca cierto porcentaje acumulado?}
  \item \textbf{Ejemplo:} estar en el percentil 80 significa estar igual o por encima de cerca del 80\% de los datos, según convención.
  \item \errorcomun{confundir percentil 80 con obtener 80 puntos.}
\end{itemize}
\end{frame}

\begin{frame}{Ejemplo de percentil con interpolación}
\textbf{Datos agrupados:} $n=50$. Se busca $P_{70}$, posición $0{,}70n=35$.
\begin{itemize}
  \item Intervalo que contiene la posición 35: [60,70).
  \item Datos del intervalo: $L=60$, $F_{ant}=30$, $f_i=10$, $A=10$.
\end{itemize}
\[
P_{70}=60+\left(\frac{35-30}{10}\right)10=65
\]
\begin{block}{Interpretación}
Aproximadamente el 70\% de los datos está por debajo de 65.
\end{block}
\end{frame}

\begin{frame}{Cuartiles}
\textbf{Definición formal:} cuantiles que dividen los datos ordenados en cuatro partes.
\begin{itemize}
  \item $Q_1$: primer cuartil, deja cerca del 25\% de los datos por debajo.
  \item $Q_2$: segundo cuartil, coincide con la mediana.
  \item $Q_3$: tercer cuartil, deja cerca del 75\% de los datos por debajo.
  \item \pregunta{¿Dónde están los cortes del 25\%, 50\% y 75\%?}
  \item \errorcomun{creer que $Q_3$ contiene el 75\% de los datos; en realidad es un valor de corte.}
\end{itemize}
\end{frame}

\begin{frame}{Deciles}
\textbf{Definición formal:} cuantiles que dividen los datos ordenados en diez partes de 10\% cada una.
\begin{itemize}
  \item $D_1$ deja cerca del 10\% de los datos por debajo.
  \item $D_5$ coincide con la mediana, bajo convenciones usuales.
  \item $D_9$ deja cerca del 90\% por debajo.
  \item \textbf{Ejemplo práctico:} si un ingreso está en $D_8$, pertenece aproximadamente al 20\% superior.
  \item \diferencia{deciles dividen en 10 partes; cuartiles en 4; percentiles en 100.}
\end{itemize}
\end{frame}

\begin{frame}{Resumen visual: cuantiles}
\centering
\begin{tikzpicture}[x=0.095cm,y=0.9cm]
  \draw[->] (0,0)--(105,0) node[right]{porcentaje acumulado};
  \foreach \x/\lab in {25/$Q_1$,50/$Q_2=Me$,75/$Q_3$}{
    \draw[thick] (\x,-0.12)--(\x,0.12) node[above]{\lab} node[below=5pt]{\x\%};
  }
  \foreach \x in {10,20,30,40,60,70,80,90}{\draw (\x,-0.07)--(\x,0.07);}
  \node[below] at (10,-0.2) {$D_1$};
  \node[below] at (90,-0.2) {$D_9$};
\end{tikzpicture}
\vspace{0.7em}
\begin{itemize}
  \item Percentiles, deciles y cuartiles son familias de cuantiles.
  \item Todos requieren datos ordenados o frecuencias acumuladas.
\end{itemize}
\end{frame}

\subsection{Medidas de dispersión}

\begin{frame}{Introducción a la dispersión}
\textbf{Definición formal:} grado en que los datos se alejan, varían o se extienden respecto de una referencia, usualmente el centro.
\begin{itemize}
  \item \textbf{Intuición:} dos grupos pueden tener el mismo promedio y comportarse de forma muy distinta.
  \item \pregunta{¿Qué tan homogéneos o heterogéneos son los datos?}
  \item \diferencia{posición indica dónde están los datos; dispersión indica qué tan separados están.}
  \item \errorcomun{comparar solo medias sin revisar variabilidad.}
\end{itemize}
\end{frame}

\begin{frame}{Rango como medida de dispersión}
\textbf{Definición formal:} diferencia entre el valor máximo y el mínimo: $R=x_{\max}-x_{\min}$.
\begin{itemize}
  \item \textbf{Intuición:} distancia entre extremos.
  \item \pregunta{¿Cuál es la extensión total de los datos?}
  \item \textbf{Ejemplo:} 12, 15, 18, 25 $\Rightarrow R=25-12=13$.
  \item \diferencia{es simple, pero solo usa dos valores; varianza y desviación usan todos los datos.}
  \item \errorcomun{concluir estabilidad solo porque el rango es pequeño sin mirar el contexto o tamaño de escala.}
\end{itemize}
\end{frame}

\begin{frame}{Varianza: concepto}
\textbf{Definición formal:} promedio de las desviaciones cuadráticas respecto de la media.
\[
\sigma^2=\frac{\sum (x_i-\mu)^2}{N}\qquad s^2=\frac{\sum (x_i-\bar{x})^2}{n-1}
\]
\begin{itemize}
  \item \textbf{Intuición:} mide qué tan lejos están los datos del promedio, penalizando más las distancias grandes.
  \item \pregunta{¿Cuánta variabilidad cuadrática hay alrededor de la media?}
  \item \diferencia{varianza poblacional divide por $N$; varianza muestral divide por $n-1$.}
  \item \errorcomun{olvidar que sus unidades están al cuadrado.}
\end{itemize}
\end{frame}

\begin{frame}{Ejemplo: varianza no agrupada}
Datos: 2, 4, 4, 6. Media: $\bar{x}=4$.
\begin{center}
\begin{tabular}{ccc}
\toprule
$x_i$ & $x_i-\bar{x}$ & $(x_i-\bar{x})^2$ \\
\midrule
2 & -2 & 4 \\
4 & 0 & 0 \\
4 & 0 & 0 \\
6 & 2 & 4 \\
\bottomrule
\end{tabular}
\end{center}
\begin{itemize}
  \item Suma de cuadrados: 8.
  \item Varianza poblacional: $\sigma^2=8/4=2$.
  \item Varianza muestral: $s^2=8/(4-1)=2{,}67$.
\end{itemize}
\end{frame}

\begin{frame}{Varianza para datos agrupados}
\textbf{Variable discreta agrupada:}
\[
s^2=\frac{\sum f_i(x_i-\bar{x})^2}{n-1}
\]
\textbf{Datos por intervalos:}
\[
s^2\approx\frac{\sum f_i(m_i-\bar{x})^2}{n-1}
\]
\begin{itemize}
  \item Para intervalos se reemplaza cada clase por su marca de clase $m_i$.
  \item \diferencia{la fórmula agrupada pondera por frecuencias; la no agrupada trabaja observación por observación.}
  \item \errorcomun{usar límites de intervalos en lugar de marcas de clase para aproximar.}
\end{itemize}
\end{frame}

\begin{frame}{Desviación estándar}
\textbf{Definición formal:} raíz cuadrada positiva de la varianza.
\[
\sigma=\sqrt{\sigma^2}\qquad s=\sqrt{s^2}
\]
\begin{itemize}
  \item \textbf{Intuición:} distancia típica de los datos respecto de la media.
  \item \pregunta{¿Cuánto se alejan normalmente los datos del promedio?}
  \item \textbf{Ejemplo:} si $s^2=2{,}67$, entonces $s=\sqrt{2{,}67}=1{,}63$.
  \item \diferencia{vuelve a las unidades originales; la varianza queda en unidades al cuadrado.}
  \item \errorcomun{interpretar una desviación estándar sin mencionar unidades o contexto.}
\end{itemize}
\end{frame}

\begin{frame}{Coeficiente de variación}
\textbf{Definición formal:} medida de dispersión relativa respecto de la media.
\[
CV=\frac{s}{\bar{x}}\quad\text{o}\quad CV\%=\frac{s}{\bar{x}}\times 100
\]
\begin{itemize}
  \item \textbf{Intuición:} variabilidad por unidad promedio.
  \item \pregunta{¿Qué tan grande es la dispersión comparada con el nivel medio?}
  \item \textbf{Ejemplo:} si $\bar{x}=50$ y $s=5$, entonces $CV=0{,}10=10\%$.
  \item \diferencia{permite comparar variables con escalas o unidades distintas.}
  \item \errorcomun{usarlo cuando la media es cero o cercana a cero.}
\end{itemize}
\end{frame}

\begin{frame}{Ejemplo comparativo: coeficiente de variación}
\begin{center}
\begin{tabular}{lccc}
\toprule
Grupo & Media & Desv. estándar & $CV\%$ \\
\midrule
A & 50 & 5 & 10\% \\
B & 200 & 20 & 10\% \\
C & 40 & 12 & 30\% \\
\bottomrule
\end{tabular}
\end{center}
\begin{itemize}
  \item A y B tienen distinta escala, pero igual dispersión relativa.
  \item C es más heterogéneo porque su desviación representa mayor proporción de la media.
  \item \textbf{Interpretación:} CV bajo sugiere mayor homogeneidad relativa.
\end{itemize}
\end{frame}

\begin{frame}{Resumen visual: dispersión}
\small
\begin{tabular}{p{0.24\textwidth}p{0.33\textwidth}p{0.31\textwidth}}
\toprule
Medida & Qué resume & Limitación o cuidado \\
\midrule
Rango & extremos & sensible a máximo y mínimo \\
Varianza & distancia cuadrática media & unidades al cuadrado \\
Desv. estándar & distancia típica en unidades originales & sensible a extremos \\
Coef. variación & dispersión relativa & no usar con media cercana a cero \\
\bottomrule
\end{tabular}
\vspace{0.7em}
\begin{block}{Idea central}
La dispersión complementa a la tendencia central: el centro sin variabilidad puede ser engañoso.
\end{block}
\end{frame}

\begin{frame}{Checklist de la Unidad 2}
\begin{enumerate}
  \item ¿Elijo el gráfico según tipo de variable y objetivo comunicativo?
  \item ¿Distingo barras, histograma, torta, ojiva, tallo-hoja y dispersión?
  \item ¿Sé cuándo usar media, mediana, moda y media ponderada?
  \item ¿Calculo medidas agrupadas usando frecuencias y marcas de clase?
  \item ¿Interpreto percentiles, cuartiles y deciles como posiciones de corte?
  \item ¿Comparo dispersión con rango, varianza, desviación estándar y CV?
\end{enumerate}
\end{frame}

\begin{frame}{Cierre integrado de la guía}
\centering
\Large
Primero se organizan los datos en tablas; luego se visualizan con gráficos y se resumen con medidas de posición y dispersión.
\vspace{0.8em}

\normalsize
Una buena interpretación combina contexto, tipo de variable, forma de la distribución y variabilidad.
\end{frame}


% -----------------------------------------------------------------------------
% Unidad 3 integrada: Números Índices
% -----------------------------------------------------------------------------
\part{ME\_3: Números Índices}

\begin{frame}{Transición a la Unidad 3}
\centering
\Large
Luego de organizar, graficar y resumir datos, el siguiente paso es comparar magnitudes en el tiempo.
\vspace{0.8em}

\normalsize
Los números índices permiten estudiar cambios relativos en precios, cantidades y valores.
\end{frame}

\begin{frame}{Cómo usar esta unidad}
\begin{itemize}
  \item Esta unidad continúa el bloque de estadística descriptiva: ahora el foco está en comparar magnitudes entre períodos.
  \item La estructura se mantiene: definición, intuición, pregunta, diferencia, error común y ejemplo.
  \item Los ejercicios típicos combinan precios, cantidades, valores, ponderaciones y elección de año base.
\end{itemize}
\end{frame}

\begin{frame}{Mapa de la Unidad 3}
\small
\begin{columns}[T]
\column{0.48\textwidth}
\begin{enumerate}
  \item Introducción a números índices
  \item Índices simples: precio, cantidad y valor
  \item Índices complejos no ponderados
  \item Índices económicos clásicos
  \item Índices ponderados
\end{enumerate}
\column{0.48\textwidth}
\begin{enumerate}\setcounter{enumi}{5}
  \item Laspeyres
  \item Paasche
  \item Relación precio--cantidad--valor
  \item Fisher
  \item Errores comunes
\end{enumerate}
\end{columns}
\end{frame}

\section{Introducción a los números índices}

\begin{frame}{¿Qué es un número índice?}
\textbf{Definición formal:} razón que compara el nivel de una variable en un período con su nivel en un período base.
\[
I_t=\frac{x_t}{x_0}
\]
\begin{itemize}
  \item \intuicion{mide cuánto representa el valor actual respecto del valor de referencia.}
  \item \pregunta{¿Cuánto cambió una magnitud en relación con el período base?}
  \item Es una razón adimensional: las unidades se cancelan.
  \item \textbf{Ejemplo:} si un precio sube de \$500 a \$650, entonces $I=650/500=1{,}30$.
  \item \errorcomun{confundir el índice $1{,}30$ con un aumento de $1{,}30\%$; el aumento es $30\%$.}
\end{itemize}
\end{frame}

\begin{frame}{Interpretación básica de un índice}
\begin{center}
\begin{tabular}{c p{0.58\textwidth} c}
\toprule
Índice & Interpretación & Variación \\
\midrule
$I=1$ & el nivel actual coincide con el período base & $0\%$ \\
$I>1$ & el nivel actual es mayor que el período base & positiva \\
$I<1$ & el nivel actual es menor que el período base & negativa \\
\bottomrule
\end{tabular}
\end{center}
\begin{block}{Conversión práctica}
Un índice puede expresarse como base 1 o base 100: $1{,}25=125$ si se usa base 100.
\end{block}
\diferencia{el índice compara niveles; la variación porcentual expresa el cambio neto respecto de la base.}
\end{frame}

\begin{frame}{Año base o período base}
\textbf{Definición formal:} período usado como referencia para comparar todos los demás períodos.
\begin{itemize}
  \item \intuicion{es el ``punto cero'' de la comparación temporal.}
  \item \pregunta{¿Respecto de qué momento estoy midiendo el cambio?}
  \item Conviene elegir un período normal, representativo y con datos confiables.
  \item Mantener una referencia fija permite comparar varios años entre sí mediante el mismo patrón.
  \item \errorcomun{cambiar el año base sin avisar y comparar índices que ya no tienen la misma referencia.}
\end{itemize}
\end{frame}

\begin{frame}{Variación porcentual a partir de un índice}
\[
\%\Delta=(I-1)\cdot 100
\]
\begin{itemize}
  \item Si $I=1{,}18$, entonces $\%\Delta=(1{,}18-1)100=18\%$.
  \item Si $I=0{,}92$, entonces $\%\Delta=(0{,}92-1)100=-8\%$.
  \item En base 100, la variación se obtiene como $I_{100}-100$.
  \item \pregunta{¿El nivel actual aumentó o disminuyó, y en qué porcentaje?}
  \item \errorcomun{decir que $I=92$ significa aumento de $92\%$; en base 100 significa caída de $8\%$.}
\end{itemize}
\end{frame}

\section{Números índices simples}

\begin{frame}{Concepto de índice simple}
\textbf{Definición formal:} índice que estudia la evolución de una sola variable o magnitud entre dos períodos.
\[
I(x)=\frac{x_t}{x_0}
\]
\begin{itemize}
  \item \intuicion{compara una cosa contra sí misma en otro momento.}
  \item \textbf{Magnitudes típicas:} precio, cantidad vendida, producción, ingreso o valor monetario.
  \item \diferencia{simple significa una variable o producto; complejo significa conjunto de productos o magnitudes.}
\end{itemize}
\end{frame}

\begin{frame}{Índice simple de precios}
\[
I(P)=\frac{P_t}{P_0}
\]
\begin{itemize}
  \item Mide la evolución temporal del precio de un bien.
  \item \pregunta{¿Cuánto cambió el precio respecto del período base?}
  \item \textbf{Ejemplo:} pan de \$1.000 a \$1.250: $I(P)=1{,}25$, aumento de $25\%$.
  \item \textbf{Interpretación económica:} aproxima inflación de ese producto específico.
  \item \errorcomun{generalizar la inflación de un producto a toda la economía.}
\end{itemize}
\end{frame}

\begin{frame}{Índice simple de cantidades}
\[
I(Q)=\frac{Q_t}{Q_0}
\]
\begin{itemize}
  \item Mide cambios en unidades físicas, producción, consumo o ventas.
  \item \pregunta{¿Se vendió o produjo más o menos que en la base?}
  \item \textbf{Ejemplo:} ventas de 80 a 100 unidades: $I(Q)=1{,}25$, aumento de $25\%$.
  \item \diferencia{un precio puede subir aunque la cantidad vendida baje; son dimensiones distintas.}
\end{itemize}
\end{frame}

\begin{frame}{Índice simple de valor}
\textbf{Valor:} monto monetario que combina precio y cantidad.
\[
V=P\cdot Q \qquad I(V)=\frac{V_t}{V_0}=\frac{P_tQ_t}{P_0Q_0}
\]
\begin{itemize}
  \item \intuicion{resume ingresos, ventas monetarias o gasto total.}
  \item Relación fundamental en índices simples:
  \[
  I(V)=I(P)\cdot I(Q)
  \]
  \item \textbf{Ejemplo:} si el precio sube $20\%$ y la cantidad baja $10\%$, entonces $I(V)=1{,}20\cdot0{,}90=1{,}08$.
  \item \errorcomun{interpretar un mayor valor como mayor cantidad vendida; puede deberse solo a precios.}
\end{itemize}
\end{frame}

\section{Números índices complejos}

\begin{frame}{Concepto de índice complejo}
\textbf{Definición formal:} índice que resume la evolución conjunta de varias magnitudes o productos.
\begin{itemize}
  \item \intuicion{pasa de mirar un producto a mirar una canasta.}
  \item \pregunta{¿Cómo cambió el conjunto completo de bienes o variables?}
  \item Se usa cuando interesa una visión global: canasta familiar, producción industrial, ventas de una empresa.
  \item \diferencia{el índice simple no requiere agregar productos; el complejo necesita un método de agregación.}
\end{itemize}
\end{frame}

\begin{frame}{Índices complejos no ponderados}
\textbf{Concepto:} todos los elementos reciben la misma importancia estadística.
\begin{itemize}
  \item \textbf{Ventaja:} cálculo rápido y fácil de explicar.
  \item \textbf{Limitación:} trata igual bienes muy distintos en relevancia económica.
  \item \textbf{Ejemplo:} dar el mismo peso a pan, arriendo y lápices puede distorsionar una canasta de consumo.
  \item \errorcomun{creer que no ponderar significa ser neutral; en realidad asigna igual peso a todo.}
\end{itemize}
\end{frame}

\begin{frame}{Media agregativa simple}
\[
I_A(P)=\frac{\sum P_t}{\sum P_0}
\]
\begin{itemize}
  \item Suma precios actuales y los compara con la suma de precios base.
  \item \pregunta{¿La suma total de magnitudes aumentó o disminuyó?}
  \item \textbf{Aplicación:} útil si las magnitudes son comparables y tienen escala similar.
  \item \errorcomun{sumar precios de bienes con unidades muy distintas sin justificar la agregación.}
\end{itemize}
\end{frame}

\begin{frame}{Método de los promedios}
\[
I_M(P)=\frac{1}{n}\sum_{i=1}^{n}\frac{P_{it}}{P_{i0}}
\]
\begin{itemize}
  \item Calcula primero índices individuales y luego los promedia.
  \item \pregunta{¿Cuál es el cambio relativo promedio de los productos?}
  \item \diferencia{la agregativa compara sumas; el método de promedios compara razones individuales.}
  \item Puede diferir mucho de la media agregativa cuando los precios base tienen escalas distintas.
\end{itemize}
\end{frame}

\begin{frame}{Comparación: agregativa simple vs promedio simple}
\small
\begin{center}
\begin{tabular}{lccccc}
\toprule
Producto & $P_0$ & $P_t$ & $P_t/P_0$ & Aporte a sumas & Aporte al promedio \\
\midrule
A & 100 & 120 & 1,20 & alto & igual \\
B & 10 & 20 & 2,00 & bajo & igual \\
\bottomrule
\end{tabular}
\end{center}
\[
I_A=\frac{120+20}{100+10}=1{,}27 \qquad
I_M=\frac{1{,}20+2{,}00}{2}=1{,}60
\]
\begin{block}{Idea clave}
Los métodos divergen cuando un producto pequeño cambia mucho o cuando las escalas iniciales son muy diferentes.
\end{block}
\end{frame}

\section{Índices económicos clásicos}

\begin{frame}{Índices económicos clásicos}
\begin{columns}[T]
\column{0.32\textwidth}
\textbf{Precios}
\begin{itemize}
  \item Inflación
  \item Poder adquisitivo
  \item Comparación temporal
\end{itemize}
\column{0.32\textwidth}
\textbf{Cantidades}
\begin{itemize}
  \item Producción
  \item Consumo
  \item Ventas físicas
\end{itemize}
\column{0.32\textwidth}
\textbf{Valor}
\begin{itemize}
  \item Ingresos
  \item Gasto total
  \item Efecto precio y cantidad
\end{itemize}
\end{columns}
\vspace{0.5em}
\errorcomun{concluir causalidad económica solo porque un índice aumentó o disminuyó.}
\end{frame}

\section{Índices compuestos ponderados}

\begin{frame}{Motivación de las ponderaciones}
\textbf{Idea central:} no todos los productos tienen la misma importancia en una canasta o en una economía.
\begin{itemize}
  \item Las ponderaciones representan participación, relevancia o peso económico.
  \item Evitan que un bien irrelevante tenga el mismo impacto que uno fundamental.
  \item \textbf{Ejemplo:} en el presupuesto familiar, el arriendo pesa más que un paquete de lápices.
  \item \pregunta{¿Qué productos deberían influir más en el índice global?}
\end{itemize}
\end{frame}

\begin{frame}{Concepto de ponderación}
\begin{itemize}
  \item \textbf{Peso estadístico:} coeficiente usado para dar importancia relativa a cada elemento.
  \item \textbf{Peso económico:} importancia medida por gasto, ingreso, producción o participación.
  \item En índices de precios, las cantidades suelen actuar como ponderaciones.
  \item En índices de cantidades, los precios suelen actuar como ponderaciones.
\end{itemize}
\begin{block}{Regla de lectura}
Precio se pondera con cantidad; cantidad se pondera con precio.
\end{block}
\end{frame}

\section{Índices de Laspeyres}

\begin{frame}{Índice de precios de Laspeyres}
\[
P_L=\frac{\sum P_tQ_0}{\sum P_0Q_0}
\]
\begin{itemize}
  \item Usa cantidades del período base como ponderaciones.
  \item \pregunta{¿Cuánto costaría hoy comprar la canasta base?}
  \item \textbf{Ventaja:} mantiene fija la canasta y facilita comparaciones.
  \item \textbf{Sesgo típico:} puede sobreestimar aumentos si los consumidores sustituyen bienes encarecidos. 
  \item \errorcomun{usar cantidades actuales en Laspeyres de precios.}
\end{itemize}
\end{frame}

\begin{frame}{Índice de cantidades de Laspeyres}
\[
Q_L=\frac{\sum Q_tP_0}{\sum Q_0P_0}
\]
\begin{itemize}
  \item Usa precios del período base como ponderaciones.
  \item \pregunta{¿Cómo cambió el volumen producido o vendido valorado a precios base?}
  \item \textbf{Aplicación:} comparar producción física evitando que los precios actuales distorsionen el cambio.
\end{itemize}
\end{frame}

\section{Índices de Paasche}

\begin{frame}{Índice de precios de Paasche}
\[
P_P=\frac{\sum P_tQ_t}{\sum P_0Q_t}
\]
\begin{itemize}
  \item Usa cantidades actuales como ponderaciones.
  \item \pregunta{¿Cuánto cuesta la canasta actual comparada a precios base?}
  \item Refleja patrones actuales de consumo o producción.
  \item \errorcomun{confundirlo con Laspeyres porque ambos comparan precios; la diferencia está en las cantidades usadas.}
\end{itemize}
\end{frame}

\begin{frame}{Índice de cantidades de Paasche}
\[
Q_P=\frac{\sum Q_tP_t}{\sum Q_0P_t}
\]
\begin{itemize}
  \item Usa precios actuales como ponderaciones.
  \item \pregunta{¿Cómo cambió el volumen usando la estructura de precios actual?}
  \item Puede incorporar cambios de importancia económica ocurridos en el período actual.
\end{itemize}
\end{frame}

\begin{frame}{Laspeyres vs Paasche}
\small
\begin{center}
\begin{tabular}{p{0.22\textwidth}p{0.33\textwidth}p{0.33\textwidth}}
\toprule
Criterio & Laspeyres & Paasche \\
\midrule
Ponderación & base & actual \\
Canasta & fija en período base & actual o móvil \\
Lectura & costo de la canasta base hoy & costo de la canasta actual vs base \\
Riesgo & puede ignorar sustitución & puede cambiar la referencia implícita \\
\bottomrule
\end{tabular}
\end{center}
\diferencia{Laspeyres responde con la estructura base; Paasche responde con la estructura actual.}
\end{frame}

\section{Relación entre índices}

\begin{frame}{Relación precio, cantidad y valor}
\[
V=\sum PQ
\]
\begin{itemize}
  \item El valor monetario cambia por dos fuerzas: precios y cantidades.
  \item En el caso simple se cumple exactamente $I(V)=I(P)I(Q)$.
  \item En índices agregados, la relación depende de que las fórmulas de precio y cantidad sean compatibles.
  \item \pregunta{¿El cambio en ventas monetarias se explica por precios, cantidades o ambos?}
\end{itemize}
\end{frame}

\begin{frame}{Identidad fundamental de valor}
\[
I_V=\frac{\sum P_tQ_t}{\sum P_0Q_0}
\]
\begin{block}{Significado económico}
Compara el valor total actual con el valor total del período base.
\end{block}
\begin{itemize}
  \item Si $I_V=1{,}40$, el valor total aumentó $40\%$.
  \item No permite separar por sí solo cuánto se debe a precio y cuánto a cantidad.
  \item \errorcomun{atribuir todo el aumento de valor a inflación sin revisar cantidades.}
\end{itemize}
\end{frame}

\section{Índices de Fisher}

\begin{frame}{Motivación del índice de Fisher}
\begin{itemize}
  \item Laspeyres usa ponderaciones base y Paasche usa ponderaciones actuales.
  \item Cada uno puede entregar resultados distintos cuando cambian los patrones de consumo o producción.
  \item Fisher busca un índice equilibrado combinando ambos.
  \item \pregunta{¿Existe una medida intermedia entre Laspeyres y Paasche?}
\end{itemize}
\end{frame}

\begin{frame}{Índice de precios de Fisher}
\[
P_F=\sqrt{P_L\cdot P_P}
\]
\begin{itemize}
  \item Es la media geométrica entre el índice de precios de Laspeyres y el de Paasche.
  \item \intuicion{equilibra la mirada de canasta base y canasta actual.}
  \item \textbf{Ventaja:} reduce la dependencia de una única estructura de ponderaciones.
  \item \errorcomun{calcular Fisher con media aritmética en vez de media geométrica.}
\end{itemize}
\end{frame}

\begin{frame}{Índice de cantidades de Fisher}
\[
Q_F=\sqrt{Q_L\cdot Q_P}
\]
\begin{itemize}
  \item Combina el índice de cantidades de Laspeyres y el de Paasche.
  \item \pregunta{¿Cuál es el cambio físico equilibrando precios base y actuales?}
  \item Es útil cuando ninguna estructura de precios parece suficiente por sí sola.
\end{itemize}
\end{frame}

\begin{frame}{Comparación final: Laspeyres, Paasche y Fisher}
\small
\begin{center}
\begin{tabular}{p{0.20\textwidth}p{0.27\textwidth}p{0.27\textwidth}p{0.16\textwidth}}
\toprule
Índice & Pondera con & Fortalezas & Cuidado \\
\midrule
Laspeyres & período base & simple, comparable, canasta fija & sustitución \\
Paasche & período actual & refleja estructura actual & base implícita móvil \\
Fisher & media geométrica & equilibrio entre ambos & requiere calcular dos índices \\
\bottomrule
\end{tabular}
\end{center}
\begin{block}{Regla práctica}
Use Laspeyres para seguimiento con canasta fija, Paasche para estructura actual y Fisher como síntesis equilibrada.
\end{block}
\end{frame}

\section{Ejemplo integrador}

\begin{frame}{Ejemplo integrador: datos}
\small
\begin{center}
\begin{tabular}{lrrrr}
\toprule
Producto & $P_0$ & $Q_0$ & $P_t$ & $Q_t$ \\
\midrule
A & 10 & 100 & 12 & 90 \\
B & 20 & 50 & 26 & 60 \\
C & 5 & 200 & 6 & 240 \\
\bottomrule
\end{tabular}
\end{center}
\begin{itemize}
  \item $\sum P_0Q_0=3000$.
  \item $\sum P_tQ_0=3700$.
  \item $\sum P_0Q_t=3300$.
  \item $\sum P_tQ_t=4080$.
\end{itemize}
\end{frame}

\begin{frame}{Ejemplo integrador: cálculo e interpretación}
\small
\[
P_L=\frac{3700}{3000}=1{,}233 \qquad
P_P=\frac{4080}{3300}=1{,}236 \qquad
P_F=\sqrt{1{,}233\cdot1{,}236}=1{,}235
\]
\[
I_V=\frac{4080}{3000}=1{,}36
\]
\begin{itemize}
  \item Los precios aumentaron aproximadamente $23{,}5\%$ según Fisher.
  \item El valor total aumentó $36\%$.
  \item Como el valor crece más que los precios, también hubo efecto de cantidades.
  \item \errorcomun{decir que el valor subió $36\%$ solo por inflación.}
\end{itemize}
\end{frame}

\section{Errores comunes}

\begin{frame}{Errores comunes en números índices}
\small
\begin{enumerate}
  \item Confundir índice con porcentaje.
  \item Interpretar $1{,}20$ como aumento de $1{,}20\%$ en vez de $20\%$.
  \item Cambiar inadvertidamente el período base.
  \item Confundir precio con valor.
  \item Confundir cantidad con valor.
  \item Aplicar Laspeyres usando ponderaciones actuales.
  \item Aplicar Paasche usando ponderaciones base.
  \item Olvidar que Fisher es una media geométrica.
  \item Interpretar causalidad económica a partir de un índice descriptivo.
\end{enumerate}
\end{frame}

\begin{frame}{Checklist de la Unidad 3}
\begin{enumerate}
  \item ¿Identifico claramente período base y período actual?
  \item ¿Distingo índice, variación porcentual e índice base 100?
  \item ¿Separo precio, cantidad y valor?
  \item ¿Reconozco si el índice es simple, complejo, ponderado o no ponderado?
  \item ¿Uso ponderaciones base para Laspeyres y actuales para Paasche?
  \item ¿Calculo Fisher como media geométrica?
  \item ¿Interpreto el resultado sin inventar causalidad?
\end{enumerate}
\end{frame}

\begin{frame}{Cierre de la unidad}
\centering
\Large
Los números índices convierten datos temporales en comparaciones relativas claras.
\vspace{0.8em}

\normalsize
La interpretación correcta depende de tres decisiones: base, ponderación y magnitud analizada.
\end{frame}

% -----------------------------------------------------------------------------
% Unidad 4 integrada: Teoría de Probabilidades y Análisis Combinatorio
% -----------------------------------------------------------------------------
\part{ME\_4: Teoría de Probabilidades y Análisis Combinatorio}

\begin{frame}{Transición a la Unidad 4}
\centering
\Large
Después de describir datos y comparar cambios, ahora modelamos la incertidumbre.
\vspace{0.8em}

\normalsize
La probabilidad conecta experimentos aleatorios, conteo y toma de decisiones bajo riesgo.
\end{frame}

\begin{frame}{Objetivos de aprendizaje de ME\_4}
\begin{columns}[T]
\column{0.50\textwidth}
\begin{block}{Al finalizar podrás}
\begin{itemize}
  \item Identificar experimentos aleatorios.
  \item Construir espacios muestrales.
  \item Operar con eventos.
  \item Contar resultados sin enumerarlos todos.
\end{itemize}
\end{block}
\column{0.45\textwidth}
\begin{exampleblock}{Preparación para}
\begin{itemize}
  \item probabilidad condicionada;
  \item variables aleatorias;
  \item inferencia estadística;
  \item modelos de confiabilidad.
\end{itemize}
\end{exampleblock}
\end{columns}
\end{frame}

\begin{frame}{Mapa visual de la unidad}
\centering
\begin{tikzpicture}[node distance=0.85cm, box/.style={rectangle, rounded corners, draw=blue!65, fill=blue!8, align=center, minimum width=2.7cm, minimum height=0.72cm}, arrow/.style={-{Latex}, thick}]
\node[box] (exp) {Experimento\\aleatorio};
\node[box, right=of exp] (omega) {Espacio\\muestral $\Omega$};
\node[box, right=of omega] (eventos) {Eventos};
\node[box, below=of omega] (alg) {Álgebra\\de eventos};
\node[box, below=of exp] (conteo) {Conteo};
\node[box, below=of eventos] (prob) {Probabilidad};
\draw[arrow] (exp) -- (omega);
\draw[arrow] (omega) -- (eventos);
\draw[arrow] (eventos) -- (alg);
\draw[arrow] (conteo) -- (prob);
\draw[arrow] (eventos) -- (prob);
\draw[arrow] (alg) -- (prob);
\end{tikzpicture}
\vspace{0.4em}
\begin{alertblock}{Idea central}
Para calcular probabilidades primero hay que saber \textbf{qué puede ocurrir} y \textbf{cuántas formas favorables existen}.
\end{alertblock}
\end{frame}

\section{Experimentos aleatorios}

\begin{frame}{¿Qué es un experimento aleatorio?}
\begin{block}{Definición formal}
Procedimiento u observación cuyo resultado exacto no puede predecirse con certeza antes de realizarlo, aunque sí se conoce el conjunto de resultados posibles.
\end{block}
\pause
\begin{itemize}
  \item \intuicion{repetir el mismo procedimiento puede producir resultados distintos.}
  \item \textbf{Interpretación práctica:} describe incertidumbre controlada.
  \item \textbf{Aplicación:} fallas de componentes, llegada de paquetes, respuesta de un modelo de IA.
  \item \errorcomun{creer que aleatorio significa ``sin reglas''; sí hay reglas, pero no certeza del resultado individual.}
\end{itemize}
\end{frame}

\begin{frame}{Determinístico vs aleatorio}
\small
\begin{center}
\begin{tabular}{p{0.30\textwidth}p{0.30\textwidth}p{0.30\textwidth}}
\toprule
Criterio & Determinístico & Aleatorio \\
\midrule
Resultado & único y predecible & incierto antes de observar \\
Ejemplo & $2+3=5$ & lanzar un dado \\
Pregunta & ¿cuál es el resultado? & ¿qué resultados pueden ocurrir? \\
Modelo & ecuación exacta & espacio muestral y probabilidades \\
\bottomrule
\end{tabular}
\end{center}
\pause
\begin{exampleblock}{Ejercicio guiado}
Clasifica: medir la vida útil de un componente electrónico. \pause Es \textbf{aleatorio}: no conocemos su duración exacta antes de observarla.
\end{exampleblock}
\end{frame}

\begin{frame}{Ejemplos de experimentos aleatorios}
\begin{columns}[T]
\column{0.48\textwidth}
\begin{itemize}
  \item Moneda: cara o sello.
  \item Dado: valores de 1 a 6.
  \item Vida útil: tiempo hasta falla.
  \item Llegada de clientes: número por minuto.
\end{itemize}
\column{0.48\textwidth}
\begin{tikzpicture}[scale=0.82]
\node[circle, draw=blue!70, fill=blue!10, minimum size=1.2cm] at (0,1.4) {Moneda};
\node[regular polygon, regular polygon sides=4, draw=green!60!black, fill=green!10, minimum size=1.1cm] at (2.4,1.4) {Dado};
\node[rectangle, rounded corners, draw=orange!80!black, fill=orange!10, minimum width=2.2cm] at (0,0) {Componente};
\node[rectangle, rounded corners, draw=purple!70, fill=purple!10, minimum width=2.2cm] at (2.4,0) {Clientes};
\end{tikzpicture}
\end{columns}
\begin{alertblock}{Resumen visual}
Experimento $\rightarrow$ resultado incierto $\rightarrow$ conjunto de posibilidades $\rightarrow$ probabilidad.
\end{alertblock}
\end{frame}

\begin{frame}{Aplicaciones reales de la aleatoriedad}
\small
\begin{tabular}{p{0.22\textwidth}p{0.33\textwidth}p{0.35\textwidth}}
\toprule
Área & Experimento aleatorio & Decisión apoyada \\
\midrule
Telecomunicaciones & pérdida de paquetes & dimensionar redes \\
Electrónica & tiempo hasta falla & diseñar mantenimiento \\
Confiabilidad & número de fallas & estimar riesgo \\
IA & predicción correcta/incorrecta & medir desempeño \\
Finanzas & retorno diario & gestionar portafolios \\
\bottomrule
\end{tabular}
\vspace{0.5em}
\observacion{La probabilidad no elimina la incertidumbre: la cuantifica.}
\end{frame}

\section{Espacio muestral}

\begin{frame}{Espacio muestral $\Omega$}
\begin{block}{Definición formal}
El espacio muestral $\Omega$ es el conjunto de todos los resultados posibles de un experimento aleatorio.
\end{block}
\pause
\begin{columns}[T]
\column{0.50\textwidth}
\begin{itemize}
  \item Moneda: $\Omega=\{C,S\}$.
  \item Dado: $\Omega=\{1,2,3,4,5,6\}$.
  \item Dos monedas: pares ordenados.
\end{itemize}
\column{0.45\textwidth}
\begin{exampleblock}{Interpretación}
Antes de hablar de probabilidad, debemos listar qué resultados son posibles.
\end{exampleblock}
\end{columns}
\end{frame}

\begin{frame}{Tipos de espacios muestrales}
\small
\begin{tabular}{p{0.24\textwidth}p{0.30\textwidth}p{0.34\textwidth}}
\toprule
Tipo & Ejemplo & Comentario \\
\midrule
Finito & dado & se puede listar completo \\
Infinito numerable & lanzar hasta obtener 1 & lista $w_1,w_2,w_3,\ldots$ \\
Infinito no numerable & vida útil continua & requiere herramientas continuas \\
\bottomrule
\end{tabular}
\vspace{0.7em}
\begin{alertblock}{Error común}
Confundir ``muchos resultados'' con infinito. Un PIN de 4 dígitos tiene muchos resultados, pero sigue siendo finito.
\end{alertblock}
\end{frame}

\begin{frame}{Construcción paso a paso: dos monedas}
\begin{columns}[T]
\column{0.45\textwidth}
\begin{enumerate}
  \item Primera moneda: $\{C,S\}$.
  \pause
  \item Segunda moneda: $\{C,S\}$.
  \pause
  \item Resultado: par ordenado.
\end{enumerate}
\[
\Omega=\{(C,C),(C,S),(S,C),(S,S)\}
\]
\column{0.50\textwidth}
\begin{center}
\begin{tabular}{c|cc}
 & C & S \\
\hline
C & $(C,C)$ & $(C,S)$ \\
S & $(S,C)$ & $(S,S)$ \\
\end{tabular}
\end{center}
\end{columns}
\begin{exampleblock}{Ejercicio 1 guiado}
¿Cuántos resultados hay al lanzar dos monedas? \pause $2\times2=4$.
\end{exampleblock}
\end{frame}

\section{Eventos}

\begin{frame}{Evento}
\begin{block}{Definición formal}
Un evento es un subconjunto del espacio muestral: $A\subseteq\Omega$.
\end{block}
\begin{columns}[T]
\column{0.48\textwidth}
\begin{itemize}
  \item \textbf{Simple:} contiene un resultado.
  \item \textbf{Compuesto:} contiene varios resultados.
  \item \textbf{Imposible:} $\varnothing$.
  \item \textbf{Seguro:} $\Omega$.
\end{itemize}
\column{0.48\textwidth}
\begin{tikzpicture}[scale=0.82]
\draw[rounded corners, fill=blue!5] (0,0) rectangle (4,2.5);
\node at (3.7,2.2) {$\Omega$};
\draw[fill=orange!30] (1.7,1.25) circle (0.8);
\node at (1.7,1.25) {$A$};
\end{tikzpicture}
\end{columns}
\end{frame}

\begin{frame}{Eventos con un dado}
\small
Sea $\Omega=\{1,2,3,4,5,6\}$.
\begin{center}
\begin{tabular}{p{0.28\textwidth}p{0.28\textwidth}p{0.32\textwidth}}
\toprule
Evento & Conjunto & Tipo \\
\midrule
Sale 4 & $\{4\}$ & simple \\
Sale par & $\{2,4,6\}$ & compuesto \\
Sale 7 & $\varnothing$ & imposible \\
Sale número de 1 a 6 & $\Omega$ & seguro \\
\bottomrule
\end{tabular}
\end{center}
\begin{exampleblock}{Ejercicio guiado}
Evento ``sale mayor que 3'': \pause $B=\{4,5,6\}$.
\end{exampleblock}
\end{frame}

\section{Álgebra de eventos}

\begin{frame}{Álgebra de eventos: operadores básicos}
\small
\begin{tabular}{p{0.20\textwidth}p{0.28\textwidth}p{0.40\textwidth}}
\toprule
Operación & Símbolo & Interpretación \\
\midrule
Unión & $A\cup B$ & ocurre $A$ o $B$ o ambos \\
Intersección & $A\cap B$ & ocurren $A$ y $B$ a la vez \\
Complemento & $A^c$ & no ocurre $A$ \\
Diferencia & $A-B$ & ocurre $A$ pero no $B$ \\
\bottomrule
\end{tabular}
\vspace{0.6em}
\observacion{Estas operaciones permiten traducir frases del lenguaje cotidiano a conjuntos.}
\end{frame}

\begin{frame}{Unión $A\cup B$}
\begin{columns}[T]
\column{0.48\textwidth}
\begin{block}{Definición}
$A\cup B=\{x:x\in A \text{ o } x\in B\}$.
\end{block}
\begin{itemize}
  \item Incluye lo que está en $A$, en $B$ o en ambos.
  \item \errorcomun{leer ``o'' como excluyente cuando en conjuntos suele ser inclusivo.}
\end{itemize}
\column{0.48\textwidth}
\begin{tikzpicture}[scale=0.9]
\draw[rounded corners] (-.2,-.2) rectangle (3.8,2.2);
\begin{scope}
\clip (1.3,1) circle (0.8);
\fill[blue!35] (-.2,-.2) rectangle (3.8,2.2);
\end{scope}
\begin{scope}
\clip (2.2,1) circle (0.8);
\fill[blue!35] (-.2,-.2) rectangle (3.8,2.2);
\end{scope}
\draw (1.3,1) circle (0.8) node[left] {$A$};
\draw (2.2,1) circle (0.8) node[right] {$B$};
\end{tikzpicture}
\end{columns}
\end{frame}

\begin{frame}{Intersección $A\cap B$}
\begin{columns}[T]
\column{0.48\textwidth}
\begin{block}{Definición}
$A\cap B=\{x:x\in A \text{ y } x\in B\}$.
\end{block}
\begin{itemize}
  \item Solo contiene resultados comunes.
  \item Si $A\cap B=\varnothing$, son mutuamente excluyentes.
\end{itemize}
\column{0.48\textwidth}
\begin{tikzpicture}[scale=0.9]
\draw[rounded corners] (-.2,-.2) rectangle (3.8,2.2);
\begin{scope}
\clip (1.3,1) circle (0.8);
\clip (2.2,1) circle (0.8);
\fill[orange!45] (-.2,-.2) rectangle (3.8,2.2);
\end{scope}
\draw (1.3,1) circle (0.8) node[left] {$A$};
\draw (2.2,1) circle (0.8) node[right] {$B$};
\end{tikzpicture}
\end{columns}
\end{frame}

\begin{frame}{Complemento y diferencia}
\small
\begin{columns}[T]
\column{0.48\textwidth}
\begin{block}{Complemento}
$A^c=\{x\in\Omega:x\notin A\}$.
\end{block}
\begin{tikzpicture}[scale=0.75]
\fill[blue!18] (0,0) rectangle (4,2.4);
\fill[white] (2,1.2) circle (0.75);
\draw (0,0) rectangle (4,2.4);
\draw (2,1.2) circle (0.75) node {$A$};
\node at (3.5,2.05) {$A^c$};
\end{tikzpicture}
\column{0.48\textwidth}
\begin{block}{Diferencia}
$A-B=A\cap B^c$.
\end{block}
\begin{tikzpicture}[scale=0.75]
\draw (0,0) rectangle (4,2.4);
\begin{scope}
\clip (1.6,1.2) circle (0.8);
\fill[green!35] (0,0) rectangle (4,2.4);
\end{scope}
\begin{scope}
\clip (2.4,1.2) circle (0.8);
\fill[white] (0,0) rectangle (4,2.4);
\end{scope}
\draw (1.6,1.2) circle (0.8) node[left] {$A$};
\draw (2.4,1.2) circle (0.8) node[right] {$B$};
\end{tikzpicture}
\end{columns}
\end{frame}

\begin{frame}{Ejemplo resuelto: dados y eventos}
Sea $\Omega=\{1,2,3,4,5,6\}$, $A=$ pares y $B=$ mayores que 3.
\[
A=\{2,4,6\},\qquad B=\{4,5,6\}
\]
\pause
\begin{columns}[T]
\column{0.50\textwidth}
\begin{align*}
A\cup B&=\{2,4,5,6\}\\
A\cap B&=\{4,6\}
\end{align*}
\column{0.45\textwidth}
\begin{align*}
A^c&=\{1,3,5\}\\
A-B&=\{2\}
\end{align*}
\end{columns}
\begin{exampleblock}{Ejercicio guiado}
Calcula $B-A$. \pause $B-A=\{5\}$.
\end{exampleblock}
\end{frame}

\section{Producto cartesiano}

\begin{frame}{Producto cartesiano}
\begin{block}{Definición formal}
$A\times B=\{(a,b):a\in A,\ b\in B\}$.
\end{block}
\begin{itemize}
  \item Forma pares ordenados.
  \item Permite construir espacios muestrales compuestos.
  \item \textbf{Interpretación:} primero ocurre un resultado de $A$ y luego uno de $B$.
\end{itemize}
\begin{alertblock}{Error común}
$(C,1)$ y $(1,C)$ no significan lo mismo si las etapas tienen roles distintos.
\end{alertblock}
\end{frame}

\begin{frame}{Moneda + dado como producto cartesiano}
\small
\[
\Omega=\{C,S\}\times\{1,2,3,4,5,6\}
\]
\begin{center}
\begin{tabular}{c|cccccc}
 & 1 & 2 & 3 & 4 & 5 & 6 \\
\hline
C & $(C,1)$ & $(C,2)$ & $(C,3)$ & $(C,4)$ & $(C,5)$ & $(C,6)$ \\
S & $(S,1)$ & $(S,2)$ & $(S,3)$ & $(S,4)$ & $(S,5)$ & $(S,6)$ \\
\end{tabular}
\end{center}
\begin{exampleblock}{Ejercicio 2 guiado}
¿Cuántos resultados hay? \pause $2\times6=12$.
\end{exampleblock}
\end{frame}

\begin{frame}{Dos ruletas: tabla visual}
\small
Cada ruleta tiene 8 resultados: $R_1=\{1,\dots,8\}$ y $R_2=\{1,\dots,8\}$.
\begin{columns}[T]
\column{0.42\textwidth}
\begin{block}{Espacio compuesto}
\[
\Omega=R_1\times R_2
\]
Cada resultado es un par $(i,j)$.
\end{block}
\column{0.52\textwidth}
\begin{center}
\begin{tikzpicture}[scale=0.38]
\foreach \i in {1,...,8}{\foreach \j in {1,...,8}{\draw[fill=blue!8] (\i,\j) rectangle +(1,1);}}
\node at (5,9.7) {Ruleta 2};
\node[rotate=90] at (-.4,5) {Ruleta 1};
\end{tikzpicture}
\end{center}
\end{columns}
\end{frame}

\section{Reglas de enumeración}

\begin{frame}{¿Por qué contar sin enumerar?}
\begin{columns}[T]
\column{0.50\textwidth}
\begin{itemize}
  \item Contraseñas.
  \item Códigos de productos.
  \item PIN bancarios.
  \item Rutas en redes.
\end{itemize}
\column{0.45\textwidth}
\begin{alertblock}{Problema}
Enumerar todo puede ser lento, incompleto o imposible en la práctica.
\end{alertblock}
\begin{exampleblock}{Idea}
Usamos reglas de conteo para obtener el tamaño del espacio muestral.
\end{exampleblock}
\end{columns}
\end{frame}

\begin{frame}{Regla de multiplicación}
\begin{theorem}[Principio multiplicativo]
Si una tarea tiene $k$ etapas con $n_1,n_2,\ldots,n_k$ opciones, entonces hay
\[
n_1\times n_2\times\cdots\times n_k
\]
formas de realizarla.
\end{theorem}
\pause
\begin{tikzpicture}[node distance=1.1cm, box/.style={rectangle, rounded corners, draw=blue!70, fill=blue!8, align=center, minimum width=2cm}, arrow/.style={-{Latex}, thick}]
\node[box] (a) {$n_1$};
\node[box, right=of a] (b) {$n_2$};
\node[box, right=of b] (c) {$\cdots$};
\node[box, right=of c] (d) {$n_k$};
\draw[arrow] (a)--(b); \draw[arrow] (b)--(c); \draw[arrow] (c)--(d);
\end{tikzpicture}
\end{frame}

\begin{frame}{Ejercicio 2 resuelto: dos ruletas}
\begin{block}{Enunciado}
Dos ruletas tienen 8 resultados posibles cada una. ¿Cuántos resultados compuestos existen?
\end{block}
\pause
\begin{enumerate}
  \item Primera ruleta: $8$ opciones.
  \pause
  \item Segunda ruleta: $8$ opciones.
  \pause
  \item Regla de multiplicación: $8\times8=64$.
\end{enumerate}
\begin{exampleblock}{Respuesta}
El espacio muestral tiene $64$ pares ordenados.
\end{exampleblock}
\end{frame}

\begin{frame}{Regla de adición}
\begin{block}{Casos mutuamente excluyentes}
Si $A\cap B=\varnothing$, entonces
\[
n(A\cup B)=n(A)+n(B).
\]
\end{block}
\pause
\begin{block}{Caso general}
Si pueden superponerse,
\[
n(A\cup B)=n(A)+n(B)-n(A\cap B).
\]
\end{block}
\begin{alertblock}{Error común}
Sumar sin restar la intersección cuenta dos veces a quienes pertenecen a ambos conjuntos.
\end{alertblock}
\end{frame}

\begin{frame}{Inclusión--exclusión}
\small
\begin{theorem}[Dos conjuntos]
\[
n(A\cup B)=n(A)+n(B)-n(A\cap B)
\]
\end{theorem}
\begin{theorem}[Tres conjuntos]
\[
\begin{aligned}
n(A\cup B\cup C)&=n(A)+n(B)+n(C)\\
&-n(A\cap B)-n(A\cap C)-n(B\cap C)\\
&+n(A\cap B\cap C).
\end{aligned}
\]
\end{theorem}
\observacion{Aplicaciones: encuestas, bases de datos, minería de datos y depuración de registros duplicados.}
\end{frame}

\section{Variaciones}

\begin{frame}{Variaciones: el orden importa}
\begin{block}{Idea}
Una variación selecciona $k$ elementos desde $n$ disponibles y el orden de selección cambia el resultado.
\end{block}
\begin{center}
\begin{tabular}{ccc}
$AB$ & $\neq$ & $BA$ \\
\end{tabular}
\end{center}
\begin{exampleblock}{Aplicación}
Códigos, rankings, turnos de atención, secuencias de comandos.
\end{exampleblock}
\end{frame}

\begin{frame}{Comparación inicial de técnicas}
\small
\begin{center}
\begin{tabular}{p{0.25\textwidth}ccc}
\toprule
Técnica & ¿elige todos? & ¿importa orden? & ¿puede repetir? \\
\midrule
Variación simple & no necesariamente & sí & no \\
Variación con repetición & no necesariamente & sí & sí \\
Permutación simple & sí & sí & no \\
Combinación simple & no necesariamente & no & no \\
\bottomrule
\end{tabular}
\end{center}
\end{frame}

\begin{frame}{Variaciones simples}
\begin{block}{Fórmula}
\[
V(n,k)=\frac{n!}{(n-k)!}
\]
\end{block}
\begin{itemize}
  \item Primera posición: $n$ opciones.
  \item Segunda: $n-1$ opciones.
  \item Continúa hasta completar $k$ posiciones.
\end{itemize}
\end{frame}

\begin{frame}{Ejercicio 3 resuelto: urna con 5 fichas}
\begin{block}{Enunciado}
Una urna tiene 5 fichas distintas. Se extraen 2 sin reemplazo y el orden importa. ¿Cuántos resultados hay?
\end{block}
\pause
\begin{align*}
V(5,2)&=\frac{5!}{(5-2)!}\\
&=\frac{5!}{3!}=5\cdot4=20.
\end{align*}
\begin{exampleblock}{Respuesta}
Hay $20$ extracciones ordenadas posibles.
\end{exampleblock}
\end{frame}

\begin{frame}{Variaciones con repetición}
\begin{block}{Fórmula}
\[
VR(n,k)=n^k
\]
\end{block}
\begin{itemize}
  \item Hay $k$ posiciones.
  \item En cada posición se mantienen las $n$ opciones.
  \item Se usa cuando repetir está permitido.
\end{itemize}
\end{frame}

\begin{frame}{Ejercicio 4 resuelto: números que terminan en 7}
\begin{block}{Enunciado}
Formar números de 4 cifras usando $\{5,6,7,8,9\}$, con repetición permitida, que terminen en 7.
\end{block}
\pause
\begin{center}
\begin{tikzpicture}[node distance=.5cm, box/.style={rectangle, draw=blue!70, fill=blue!8, minimum width=1.2cm, minimum height=.9cm}]
\node[box] (a) {5};
\node[box, right=of a] (b) {5};
\node[box, right=of b] (c) {5};
\node[box, right=of c, fill=orange!25] (d) {1};
\node[below=.1cm of a] {mil};
\node[below=.1cm of b] {centena};
\node[below=.1cm of c] {decena};
\node[below=.1cm of d] {unidad=7};
\end{tikzpicture}
\end{center}
\[
5\times5\times5\times1=125
\]
\begin{alertblock}{Cuidado}
La restricción fija la última posición; no se multiplica por 5 allí.
\end{alertblock}
\end{frame}

\section{Permutaciones}

\begin{frame}{Permutaciones: ordenar todos los elementos}
\begin{block}{Definición}
Una permutación ordena todos los elementos disponibles.
\end{block}
\begin{columns}[T]
\column{0.45\textwidth}
\[
P(n)=n!
\]
\column{0.50\textwidth}
\begin{exampleblock}{Intuición}
Si todos participan, el problema no es seleccionar, sino ordenar.
\end{exampleblock}
\end{columns}
\end{frame}

\begin{frame}{Ejercicio 5 resuelto: letras A,V,E}
\begin{block}{Enunciado}
¿Cuántas formas hay de ordenar las letras A, V y E?
\end{block}
\[
P(3)=3!=6
\]
\begin{center}
\begin{tabular}{cccccc}
AVE & AEV & VAE & VEA & EAV & EVA
\end{tabular}
\end{center}
\begin{exampleblock}{Ejercicio guiado}
¿Por qué AVE y EVA son distintos? \pause Porque el orden cambia la palabra formada.
\end{exampleblock}
\end{frame}

\begin{frame}{Ejercicio 6 resuelto: ordenar 5,6,7,8,9}
\begin{block}{Enunciado}
Ordenar los cinco símbolos $5,6,7,8,9$ sin repetir.
\end{block}
\pause
\[
P(5)=5!=5\cdot4\cdot3\cdot2\cdot1=120
\]
\begin{exampleblock}{Interpretación}
Existen $120$ números distintos de 5 cifras usando todos los símbolos una vez.
\end{exampleblock}
\end{frame}

\begin{frame}{Ejercicio 7: permutación con restricción}
\begin{block}{Enunciado}
Ordenar $5,6,7,8,9$ con el número 9 fijo al inicio.
\end{block}
\pause
\begin{center}
\begin{tikzpicture}[node distance=.5cm, box/.style={rectangle, draw=blue!70, fill=blue!8, minimum width=1.15cm, minimum height=.85cm}]
\node[box, fill=orange!25] (a) {9};
\node[box, right=of a] (b) {4};
\node[box, right=of b] (c) {3};
\node[box, right=of c] (d) {2};
\node[box, right=of d] (e) {1};
\end{tikzpicture}
\end{center}
\[
1\times4\times3\times2\times1=4!=24
\]
\observacion{Fijar un elemento reduce el problema a ordenar los elementos restantes.}
\end{frame}

\begin{frame}{Permutaciones circulares}
\begin{block}{Fórmula}
\[
PC(n)=(n-1)!
\]
\end{block}
\begin{columns}[T]
\column{0.45\textwidth}
\begin{itemize}
  \item Mesa redonda.
  \item Sensores alrededor de una estructura.
  \item Nodos en un anillo de red.
\end{itemize}
\column{0.50\textwidth}
\begin{tikzpicture}[scale=.75]
\foreach \a/\n in {90/A,18/B,306/C,234/D,162/E}{\node[circle, draw=blue!70, fill=blue!10] at (\a:1.4) {\n};}
\draw[dashed] (0,0) circle (1.4);
\end{tikzpicture}
\end{columns}
\end{frame}

\begin{frame}{Permutaciones con repetición}
\begin{block}{Fórmula}
Si hay $n$ elementos con repeticiones $n_1,n_2,\ldots,n_k$:
\[
P=\frac{n!}{n_1!n_2!\cdots n_k!}
\]
\end{block}
\begin{alertblock}{Error común}
Usar $n!$ cuando hay símbolos repetidos sobrecuenta ordenamientos indistinguibles.
\end{alertblock}
\end{frame}

\begin{frame}{Ejercicio 8 resuelto: x,x,x,y,y}
\begin{block}{Enunciado}
¿Cuántos ordenamientos distintos tiene la secuencia con tres $x$ y dos $y$?
\end{block}
\pause
\begin{align*}
n&=5,\qquad n_x=3,\qquad n_y=2\\
P&=\frac{5!}{3!2!}=\frac{120}{6\cdot2}=10.
\end{align*}
\begin{exampleblock}{Respuesta}
Hay $10$ secuencias distintas.
\end{exampleblock}
\end{frame}

\section{Combinaciones}

\begin{frame}{Combinaciones: el orden no importa}
\begin{block}{Definición}
Una combinación selecciona elementos sin considerar el orden.
\end{block}
\begin{center}
$\{A,B\}=\{B,A\}$
\end{center}
\begin{exampleblock}{Aplicaciones}
Equipos de trabajo, comités, selección de estudiantes, muestras sin orden.
\end{exampleblock}
\end{frame}

\begin{frame}{Combinaciones simples}
\begin{block}{Fórmula}
\[
C(n,k)=\binom{n}{k}=\frac{n!}{k!(n-k)!}
\]
\end{block}
\begin{itemize}
  \item Primero se cuentan selecciones ordenadas.
  \item Luego se divide por las $k!$ formas de ordenar cada grupo.
\end{itemize}
\begin{exampleblock}{Ejercicio guiado}
Elegir 2 estudiantes entre 5: \pause $\binom{5}{2}=10$.
\end{exampleblock}
\end{frame}

\begin{frame}{Combinaciones con repetición}
\begin{block}{Fórmula}
\[
CR(n,k)=\binom{n+k-1}{k}=\frac{(n+k-1)!}{k!(n-1)!}
\]
\end{block}
\begin{itemize}
  \item Se permite escoger el mismo tipo varias veces.
  \item El orden no importa.
  \item Ejemplos: productos repetidos, distribución de recursos, pedidos con sabores repetidos.
\end{itemize}
\end{frame}

\section{Probabilidad}

\begin{frame}{¿Qué es probabilidad?}
\begin{block}{Definición conceptual}
Medida numérica entre 0 y 1 que cuantifica qué tan plausible es un evento.
\end{block}
\begin{columns}[T]
\column{0.48\textwidth}
\begin{itemize}
  \item $0$: imposible.
  \item $1$: seguro.
  \item Valores intermedios: incertidumbre.
\end{itemize}
\column{0.48\textwidth}
\begin{tikzpicture}
\draw[very thick, -{Latex}] (0,0)--(4,0);
\foreach \x/\l in {0/0,2/0.5,4/1}{\draw (\x,.1)--(\x,-.1) node[below]{\l};}
\node[above] at (0,0) {imposible};
\node[above] at (4,0) {seguro};
\end{tikzpicture}
\end{columns}
\end{frame}

\begin{frame}{Interpretaciones de probabilidad}
\small
\begin{tabular}{p{0.28\textwidth}p{0.32\textwidth}p{0.30\textwidth}}
\toprule
Enfoque & Idea & Ejemplo \\
\midrule
Clásico & casos favorables / posibles & dado justo \\
Frecuentista & proporción a largo plazo & tasa de fallas observada \\
Subjetivo & grado de creencia informado & riesgo financiero \\
\bottomrule
\end{tabular}
\vspace{0.6em}
\observacion{En esta unidad usamos principalmente probabilidad clásica en espacios finitos equiprobables.}
\end{frame}

\begin{frame}{Axiomas de Kolmogorov}
\begin{theorem}[Axiomas]
Para eventos en $\Omega$:
\begin{enumerate}
  \item $P(A)\ge 0$.
  \item $P(\Omega)=1$.
  \item Si $A\cap B=\varnothing$, entonces $P(A\cup B)=P(A)+P(B)$.
\end{enumerate}
\end{theorem}
\pause
\begin{exampleblock}{Interpretación física}
La probabilidad total disponible es 1 y se reparte entre resultados o eventos sin asignar masa negativa.
\end{exampleblock}
\end{frame}

\begin{frame}{Teoremas derivados útiles}
\small
\begin{center}
\begin{tabular}{p{0.32\textwidth}p{0.46\textwidth}}
\toprule
Resultado & Fórmula \\
\midrule
Evento imposible & $P(\varnothing)=0$ \\
Complemento & $P(A^c)=1-P(A)$ \\
Inclusión & si $A\subseteq B$, entonces $P(A)\le P(B)$ \\
Adición general & $P(A\cup B)=P(A)+P(B)-P(A\cap B)$ \\
\bottomrule
\end{tabular}
\end{center}
\begin{alertblock}{Error común}
Aplicar $P(A\cup B)=P(A)+P(B)$ sin verificar si los eventos son excluyentes.
\end{alertblock}
\end{frame}

\begin{frame}{Probabilidad clásica}
\begin{block}{Fórmula}
Si todos los resultados son equiprobables,
\[
P(A)=\frac{\text{casos favorables}}{\text{casos posibles}}=\frac{n(A)}{n(\Omega)}.
\]
\end{block}
\begin{enumerate}
  \item Construir $\Omega$.
  \item Construir el evento $A$.
  \item Contar casos favorables y posibles.
  \item Dividir e interpretar.
\end{enumerate}
\end{frame}

\begin{frame}{Ejercicio 9 resuelto: obtener un 2}
\begin{block}{Enunciado}
Lanzar un dado justo. ¿Cuál es la probabilidad de obtener un 2?
\end{block}
\pause
\[
\Omega=\{1,2,3,4,5,6\},\qquad A=\{2\}
\]
\[
P(A)=\frac{n(A)}{n(\Omega)}=\frac{1}{6}
\]
\begin{exampleblock}{Respuesta}
La probabilidad es $1/6\approx0{,}167$.
\end{exampleblock}
\end{frame}

\begin{frame}{Ejercicio 10 resuelto: obtener a lo más 2}
\begin{block}{Enunciado}
Lanzar un dado justo. ¿Cuál es la probabilidad de obtener a lo más 2?
\end{block}
\pause
\begin{enumerate}
  \item ``A lo más 2'' significa $\le 2$.
  \item Evento: $A=\{1,2\}$.
  \item Casos favorables: $2$.
  \item Casos posibles: $6$.
\end{enumerate}
\[
P(A)=\frac{2}{6}=\frac{1}{3}
\]
\end{frame}

\begin{frame}{Espacios infinitos numerables}
\begin{block}{Ejemplo}
Lanzar un dado hasta obtener un 1.
\end{block}
\begin{itemize}
  \item $w_1$: sale 1 en el primer lanzamiento.
  \item $w_2$: no sale 1 y luego sale 1.
  \item $w_3$: no sale 1 dos veces y luego sale 1.
\end{itemize}
\[
\Omega=\{w_1,w_2,w_3,\ldots\}
\]
\observacion{Es infinito porque no existe un número máximo garantizado de lanzamientos.}
\end{frame}

\begin{frame}{Sumas infinitas de probabilidades}
\begin{block}{Idea}
En espacios numerables, la probabilidad de un evento puede calcularse sumando probabilidades de resultados individuales.
\end{block}
\[
P(A)=\sum_{w_i\in A}P(w_i)
\]
\begin{columns}[T]
\column{0.32\textwidth}Distribución geométrica.
\column{0.32\textwidth}Cadenas de Markov.
\column{0.32\textwidth}Teoría de colas.
\end{columns}
\begin{alertblock}{Nivel de esta unidad}
Solo se introduce la idea; el tratamiento formal se desarrolla más adelante.
\end{alertblock}
\end{frame}

\section{Cierre de la unidad}

\begin{frame}{Mapa conceptual completo}
\centering
\begin{tikzpicture}[node distance=.8cm, box/.style={rectangle, rounded corners, draw=blue!65, fill=blue!8, align=center, minimum width=2.4cm, minimum height=.65cm}, arrow/.style={-{Latex}, thick}]
\node[box] (a) {Experimento};
\node[box, right=of a] (b) {$\Omega$};
\node[box, right=of b] (c) {Eventos};
\node[box, below=of b] (d) {Conteo};
\node[box, below=of c] (e) {Probabilidad};
\node[box, below=of d] (f) {Inferencia};
\draw[arrow] (a)--(b); \draw[arrow] (b)--(c); \draw[arrow] (c)--(e); \draw[arrow] (d)--(e); \draw[arrow] (e)--(f);
\end{tikzpicture}
\end{frame}

\begin{frame}{Resumen de fórmulas de conteo}
\small
\begin{center}
\begin{tabular}{p{0.25\textwidth}p{0.20\textwidth}p{0.18\textwidth}p{0.25\textwidth}}
\toprule
Técnica & Orden & Repetición & Fórmula \\
\midrule
Regla multiplicación & por etapas & depende & $n_1\cdots n_k$ \\
Variación simple & sí & no & $\frac{n!}{(n-k)!}$ \\
Variación con repetición & sí & sí & $n^k$ \\
Permutación simple & sí & no & $n!$ \\
Permutación circular & sí & no & $(n-1)!$ \\
Permutación con repetición & sí & sí & $\frac{n!}{n_1!\cdots n_k!}$ \\
Combinación simple & no & no & $\binom{n}{k}$ \\
Combinación con repetición & no & sí & $\binom{n+k-1}{k}$ \\
\bottomrule
\end{tabular}
\end{center}
\end{frame}

\begin{frame}{Checklist de decisión: ¿qué fórmula uso?}
\begin{enumerate}
  \item ¿El experimento ocurre por etapas? \pause Use multiplicación.
  \item ¿Estoy ordenando todos los elementos? \pause Permutación.
  \item ¿Estoy eligiendo solo algunos? \pause Variación o combinación.
  \item ¿Importa el orden? \pause Sí: variación. No: combinación.
  \item ¿Se permite repetición? \pause Use la versión con repetición.
\end{enumerate}
\end{frame}

\begin{frame}{Ejercicios integradores}
\small
\begin{enumerate}
  \item Construya $\Omega$ para moneda + dado y calcule $P(\text{cara y número par})$.
  \item Un PIN tiene 4 dígitos y permite repetición. ¿Cuántos PIN existen?
  \item Con las letras A, B, C, D, ¿cuántas claves de 2 letras sin repetición existen?
  \item De 6 estudiantes se elige un comité de 3. ¿Importa el orden?
  \item En un dado, sea $A=$ par y $B=$ mayor que 3. Calcule $P(A\cup B)$.
\end{enumerate}
\begin{exampleblock}{Meta}
Mezclar construcción de eventos, conteo y probabilidad clásica.
\end{exampleblock}
\end{frame}

\begin{frame}{Cierre de ME\_4}
\centering
\Large
La probabilidad comienza con una buena descripción del experimento y de sus eventos.
\vspace{0.8em}

\normalsize
Contar bien es el puente entre el espacio muestral y el cálculo de probabilidades.
\end{frame}

% -----------------------------------------------------------------------------
% Unidad integrada: Variables Aleatorias y Modelos Estocásticos
% -----------------------------------------------------------------------------
\part{Unidad 3: Variables Aleatorias y Modelos Estocásticos}

\begin{frame}{Transición: de eventos a variables aleatorias}
\centering
\Large
La probabilidad describe eventos; las variables aleatorias convierten resultados en números.
\vspace{0.7em}

\normalsize
Así podemos calcular distribuciones, promedios, dispersión y construir modelos estocásticos.
\end{frame}

\begin{frame}{Objetivo pedagógico de la unidad}
\begin{columns}[T]
\column{0.48\textwidth}
\begin{block}{Idea central}
Transformar resultados de un experimento aleatorio en valores numéricos.
\end{block}
\begin{itemize}
  \item Definir variables aleatorias.
  \item Distinguir discretas y continuas.
  \item Estudiar distribuciones.
\end{itemize}
\column{0.48\textwidth}
\begin{exampleblock}{Herramientas}
\begin{itemize}
  \item función de probabilidad;
  \item densidad;
  \item distribución acumulada;
  \item esperanza y varianza.
\end{itemize}
\end{exampleblock}
\end{columns}
\end{frame}

\begin{frame}{Mapa conceptual de la unidad}
\centering
\begin{tikzpicture}[node distance=.75cm, box/.style={rectangle, rounded corners, draw=blue!65, fill=blue!8, align=center, minimum width=2.55cm, minimum height=.62cm}, arrow/.style={-{Latex}, thick}]
\node[box] (m) {Modelo\\estocástico};
\node[box, below=of m] (omega) {Espacio\\muestral $\Omega$};
\node[box, below=of omega] (x) {Variable\\aleatoria $X$};
\node[box, below left=of x] (dist) {Distribución};
\node[box, below right=of x] (f) {FDA $F(x)$};
\node[box, below=of dist] (e) {Esperanza};
\node[box, below=of f] (v) {Varianza};
\draw[arrow] (m)--(omega); \draw[arrow] (omega)--(x); \draw[arrow] (x)--(dist); \draw[arrow] (x)--(f); \draw[arrow] (dist)--(e); \draw[arrow] (f)--(v);
\end{tikzpicture}
\end{frame}

\section{Modelos estocásticos}

\begin{frame}{¿Qué es un modelo estocástico?}
\begin{block}{Definición formal}
Representación matemática de un fenómeno cuyo comportamiento incluye incertidumbre y se describe mediante probabilidades.
\end{block}
\pause
\begin{itemize}
  \item \intuicion{no predice un único resultado; describe un conjunto de resultados posibles y sus probabilidades.}
  \item \textbf{Interpretación práctica:} cuantifica riesgo, variabilidad y escenarios.
  \item \errorcomun{pensar que un modelo estocástico es menos riguroso por no entregar una respuesta única.}
\end{itemize}
\end{frame}

\begin{frame}{Modelo determinístico vs probabilístico}
\small
\begin{center}
\begin{tabular}{p{0.27\textwidth}p{0.30\textwidth}p{0.30\textwidth}}
\toprule
Criterio & Determinístico & Probabilístico \\
\midrule
Salida & valor único & distribución de resultados \\
Incertidumbre & no explícita & parte del modelo \\
Pregunta & ¿cuánto será? & ¿con qué probabilidad ocurre? \\
Ejemplo & $V=IR$ ideal & falla de un componente \\
\bottomrule
\end{tabular}
\end{center}
\begin{alertblock}{Diferencia clave}
El determinístico transforma entradas en una salida; el estocástico transforma entradas en una distribución.
\end{alertblock}
\end{frame}

\begin{frame}{¿Por qué aparecen las variables aleatorias?}
\begin{columns}[T]
\column{0.48\textwidth}
\begin{block}{Problema}
Los resultados pueden ser cualitativos o estructurados: cara, sello, pares, rutas, fallas.
\end{block}
\pause
\begin{exampleblock}{Solución}
Asignar un número a cada resultado para poder calcular.
\end{exampleblock}
\column{0.48\textwidth}
\begin{tikzpicture}[node distance=.65cm, box/.style={rectangle, rounded corners, draw=blue!65, fill=blue!8, align=center, minimum width=2.4cm}, arrow/.style={-{Latex}, thick}]
\node[box] (r) {Resultados};
\node[box, below=of r] (x) {$X$};
\node[box, below=of x] (n) {Números};
\draw[arrow] (r)--(x); \draw[arrow] (x)--(n);
\end{tikzpicture}
\end{columns}
\end{frame}

\begin{frame}{Aplicaciones de modelos estocásticos}
\small
\begin{tabular}{p{0.23\textwidth}p{0.35\textwidth}p{0.30\textwidth}}
\toprule
Área & Variable aleatoria típica & Decisión \\
\midrule
Confiabilidad & tiempo hasta falla & mantenimiento \\
Electrónica & ruido o voltaje & tolerancias \\
Minería & ley del mineral & planificación \\
Economía & demanda o retorno & riesgo \\
Manufactura & piezas defectuosas & control de calidad \\
IA & error de predicción & evaluación \\
Telecomunicaciones & paquetes perdidos & capacidad \\
\bottomrule
\end{tabular}
\end{frame}

\section{Variables aleatorias}

\begin{frame}{Variable aleatoria: definición formal}
\begin{block}{Definición}
Una variable aleatoria es una función que asigna un número real a cada resultado del espacio muestral:
\[
X:\Omega\to\mathbb{R}.
\]
\end{block}
\pause
\begin{itemize}
  \item \textbf{Dominio:} espacio muestral $\Omega$.
  \item \textbf{Recorrido o rango:} valores numéricos que puede tomar $X$.
  \item \textbf{Interpretación:} $X(\omega)$ es el número asociado al resultado $\omega$.
\end{itemize}
\end{frame}

\begin{frame}{Transformar resultados en números}
\centering
\begin{tikzpicture}[node distance=1cm, leftset/.style={ellipse, draw=blue!70, fill=blue!8, minimum width=2.2cm, minimum height=3cm}, rightset/.style={ellipse, draw=green!60!black, fill=green!8, minimum width=2.2cm, minimum height=3cm}, arrow/.style={-{Latex}, thick}]
\node[leftset] (omega) at (0,0) {};
\node at (0,1.75) {$\Omega$};
\node at (0,.8) {$C$}; \node at (0,0) {$S$};
\node[rightset] (r) at (4,0) {};
\node at (4,1.75) {Rango};
\node at (4,.8) {$1$}; \node at (4,0) {$0$};
\draw[arrow] (.45,.8)--(3.55,.8) node[midway, above] {$X$};
\draw[arrow] (.45,0)--(3.55,0);
\end{tikzpicture}
\vspace{0.4em}
\begin{exampleblock}{Ejemplo}
Si $X=$ número de caras en una moneda, entonces $X(C)=1$ y $X(S)=0$.
\end{exampleblock}
\end{frame}

\begin{frame}{Dominio, variable aleatoria y rango}
\begin{center}
\begin{tikzpicture}[node distance=.9cm, box/.style={rectangle, rounded corners, draw=blue!65, fill=blue!8, align=center, minimum width=3.2cm, minimum height=.7cm}, arrow/.style={-{Latex}, thick}]
\node[box] (omega) {Espacio muestral\\$\Omega$};
\node[box, below=of omega] (x) {Variable aleatoria\\$X:\Omega\to\mathbb{R}$};
\node[box, below=of x] (range) {Rango\\$R_X$};
\draw[arrow] (omega)--(x); \draw[arrow] (x)--(range);
\end{tikzpicture}
\end{center}
\begin{alertblock}{Error común}
Confundir $\Omega$ con el rango: $\Omega$ contiene resultados; el rango contiene números asignados por $X$.
\end{alertblock}
\end{frame}

\begin{frame}{Ejemplo: dado}
\begin{block}{Experimento}
Lanzar un dado y definir $X=$ resultado obtenido.
\end{block}
\begin{columns}[T]
\column{0.45\textwidth}
\[
\Omega=\{1,2,3,4,5,6\}
\]
\[
R_X=\{1,2,3,4,5,6\}
\]
\column{0.48\textwidth}
\begin{exampleblock}{Interpretación}
Aquí el resultado ya es numérico, por eso $X(\omega)=\omega$.
\end{exampleblock}
\end{columns}
\end{frame}

\begin{frame}{Ejemplo: varias monedas}
\small
Lanzar tres monedas y definir $X=$ número de caras.
\begin{center}
\begin{tabular}{c c}
\toprule
Resultado $\omega$ & $X(\omega)$ \\
\midrule
CCC & 3 \\
CCS, CSC, SCC & 2 \\
CSS, SCS, SSC & 1 \\
SSS & 0 \\
\bottomrule
\end{tabular}
\end{center}
\[
R_X=\{0,1,2,3\}
\]
\end{frame}

\begin{frame}{Ejemplos de ingeniería e industria}
\small
\begin{tabular}{p{0.30\textwidth}p{0.32\textwidth}p{0.25\textwidth}}
\toprule
Experimento & Variable aleatoria & Tipo esperado \\
\midrule
Sensores en línea & número de alertas por hora & discreta \\
Medición industrial & diámetro de una pieza & continua \\
Componente electrónico & tiempo hasta falla & continua \\
Control de calidad & piezas defectuosas en lote & discreta \\
\bottomrule
\end{tabular}
\observacion{La elección de $X$ depende de la pregunta que se quiere responder.}
\end{frame}

\section{Tipos de variables aleatorias}

\begin{frame}{Variables aleatorias discretas}
\begin{block}{Definición}
Una variable aleatoria es discreta si su rango es finito o infinito numerable.
\end{block}
\begin{itemize}
  \item Se puede listar: $0,1,2,\ldots$.
  \item Sus probabilidades se asignan a puntos.
  \item Aparece al contar ocurrencias.
\end{itemize}
\begin{exampleblock}{Ejemplos}
Número de caras, clientes, fallas o piezas defectuosas.
\end{exampleblock}
\end{frame}

\begin{frame}{Variables aleatorias continuas}
\begin{block}{Definición}
Una variable aleatoria es continua si puede tomar valores en intervalos de la recta real.
\end{block}
\begin{itemize}
  \item No se enumeran punto por punto.
  \item Las probabilidades se calculan sobre intervalos.
  \item Aparece al medir magnitudes físicas.
\end{itemize}
\begin{exampleblock}{Ejemplos}
Temperatura, voltaje, corriente, tiempo, longitud, masa y resistencia eléctrica.
\end{exampleblock}
\end{frame}

\begin{frame}{Discreta vs continua}
\small
\begin{center}
\begin{tabular}{p{0.27\textwidth}p{0.30\textwidth}p{0.30\textwidth}}
\toprule
Criterio & Discreta & Continua \\
\midrule
Origen & conteo & medición \\
Valores & puntos separados & intervalos \\
Probabilidad puntual & puede ser positiva & es cero \\
Herramienta & función de probabilidad & densidad \\
Gráfico típico & barras & curva/área \\
\bottomrule
\end{tabular}
\end{center}
\end{frame}

\begin{frame}{Gráfico comparativo}
\begin{columns}[T]
\column{0.48\textwidth}
\centering
\begin{tikzpicture}
\begin{axis}[width=5cm,height=3.6cm,ybar,bar width=8pt, ymin=0,ymax=.45, xtick={0,1,2,3}, xlabel={$x$}, ylabel={$f(x)$}, title={Discreta}]
\addplot coordinates {(0,.125) (1,.375) (2,.375) (3,.125)};
\end{axis}
\end{tikzpicture}
\column{0.48\textwidth}
\centering
\begin{tikzpicture}
\begin{axis}[width=5cm,height=3.6cm, ymin=0,ymax=.45, xmin=-3,xmax=3, xlabel={$x$}, ylabel={$\rho(x)$}, title={Continua}]
\addplot[domain=-3:3, samples=80, blue, thick] {1/sqrt(2*pi)*exp(-x^2/2)};
\end{axis}
\end{tikzpicture}
\end{columns}
\end{frame}

\section{Función de probabilidad}

\begin{frame}{Función de probabilidad}
\begin{block}{Definición}
Para una variable aleatoria discreta,
\[
f(x)=P(X=x).
\]
\end{block}
\begin{columns}[T]
\column{0.48\textwidth}
\begin{itemize}
  \item $f(x)\ge0$.
  \item $\sum_x f(x)=1$.
  \item Entrega probabilidades puntuales.
\end{itemize}
\column{0.48\textwidth}
\begin{alertblock}{Error común}
Olvidar verificar que la suma de probabilidades sea 1.
\end{alertblock}
\end{columns}
\end{frame}

\begin{frame}{Ejemplo completo: tres monedas}
\small
$X=$ número de caras al lanzar tres monedas.
\begin{center}
\begin{tabular}{c c c}
\toprule
$x$ & casos & $f(x)=P(X=x)$ \\
\midrule
0 & SSS & $1/8$ \\
1 & CSS,SCS,SSC & $3/8$ \\
2 & CCS,CSC,SCC & $3/8$ \\
3 & CCC & $1/8$ \\
\bottomrule
\end{tabular}
\end{center}
\[
\sum_x f(x)=\frac18+\frac38+\frac38+\frac18=1
\]
\end{frame}

\begin{frame}{Gráfico de función de probabilidad}
\centering
\begin{tikzpicture}
\begin{axis}[width=9cm,height=5cm,ybar,bar width=12pt, ymin=0,ymax=.45, xtick={0,1,2,3}, xlabel={$x$ número de caras}, ylabel={$f(x)$}, nodes near coords]
\addplot coordinates {(0,.125) (1,.375) (2,.375) (3,.125)};
\end{axis}
\end{tikzpicture}
\end{frame}

\begin{frame}{Calcular probabilidades con $f(x)$}
\begin{block}{Pregunta}
Para tres monedas, calcular $P(X\ge2)$.
\end{block}
\pause
\begin{align*}
P(X\ge2)&=P(X=2)+P(X=3)\\
&=f(2)+f(3)\\
&=\frac38+\frac18=\frac48=\frac12.
\end{align*}
\begin{exampleblock}{Interpretación}
La probabilidad de obtener al menos dos caras es $0{,}5$.
\end{exampleblock}
\end{frame}

\section{Función de densidad}

\begin{frame}{Función de densidad}
\begin{block}{Definición}
Para una variable aleatoria continua, la densidad $\rho(x)$ permite calcular probabilidades mediante áreas.
\end{block}
\[
P(a<X<b)=\int_a^b \rho(x)\,dx
\]
\begin{itemize}
  \item $\rho(x)\ge0$.
  \item $\int_{-\infty}^{\infty}\rho(x)\,dx=1$.
  \item $P(X=x)=0$ para puntos individuales.
\end{itemize}
\end{frame}

\begin{frame}{Una densidad no es una probabilidad}
\begin{alertblock}{Cuidado central}
El valor $\rho(x)$ mide concentración relativa, no probabilidad puntual.
\end{alertblock}
\begin{columns}[T]
\column{0.48\textwidth}
\begin{itemize}
  \item Probabilidad = área.
  \item La densidad puede ser mayor que 1.
  \item Un punto no tiene ancho, por eso su área es 0.
\end{itemize}
\column{0.48\textwidth}
\begin{tikzpicture}
\begin{axis}[width=5cm,height=3.6cm, ymin=0,xmin=0,xmax=1, xlabel={$x$}, ylabel={$\rho(x)$}]
\addplot[domain=0:1, samples=2, blue, thick] {1};
\addplot[domain=.2:.6, draw=none, fill=blue!25] {1} \closedcycle;
\end{axis}
\end{tikzpicture}
\end{columns}
\end{frame}

\begin{frame}{Ejemplo continuo uniforme}
\begin{block}{Modelo}
$X\sim U(0,10)$ con densidad $\rho(x)=1/10$ para $0\le x\le10$.
\end{block}
\pause
\[
P(2<X<5)=\int_2^5 \frac{1}{10}\,dx=\frac{3}{10}=0{,}3.
\]
\begin{exampleblock}{Interpretación}
La probabilidad corresponde al área de un rectángulo de base 3 y altura 0,1.
\end{exampleblock}
\end{frame}

\begin{frame}{Área bajo la curva}
\centering
\begin{tikzpicture}
\begin{axis}[width=9cm,height=5cm, ymin=0,ymax=.45, xmin=-3,xmax=3, xlabel={$x$}, ylabel={$\rho(x)$}]
\addplot[domain=-3:3, samples=100, blue, thick] {1/sqrt(2*pi)*exp(-x^2/2)};
\addplot[domain=-1:1, samples=80, draw=none, fill=blue!25] {1/sqrt(2*pi)*exp(-x^2/2)} \closedcycle;
\draw[dashed] (axis cs:-1,0)--(axis cs:-1,.25);
\draw[dashed] (axis cs:1,0)--(axis cs:1,.25);
\end{axis}
\end{tikzpicture}
\begin{block}{Lectura}
El área sombreada representa $P(-1<X<1)$.
\end{block}
\end{frame}

\section{Función de distribución acumulada}

\begin{frame}{Función de distribución acumulada}
\begin{block}{Definición}
La función de distribución acumulada es
\[
F(x)=P(X\le x).
\]
\end{block}
\begin{itemize}
  \item Acumula probabilidad desde la izquierda.
  \item Siempre toma valores entre 0 y 1.
  \item Resume completamente la distribución de $X$.
\end{itemize}
\end{frame}

\begin{frame}{FDA discreta: función escalonada}
\centering
\begin{tikzpicture}
\begin{axis}[width=9cm,height=5cm, ymin=0,ymax=1.1, xmin=-.5,xmax=3.5, xtick={0,1,2,3}, xlabel={$x$}, ylabel={$F(x)$}]
\addplot[const plot mark right, thick, blue] coordinates {(-.5,0) (0,.125) (1,.5) (2,.875) (3,1) (3.5,1)};
\end{axis}
\end{tikzpicture}
\begin{alertblock}{Interpretación}
Los saltos ocurren en los valores posibles de la variable discreta.
\end{alertblock}
\end{frame}

\begin{frame}{FDA continua: acumulación de área}
\[
F(x)=\int_{-\infty}^{x}\rho(t)\,dt
\]
\begin{columns}[T]
\column{0.48\textwidth}
\begin{itemize}
  \item Crece suavemente si la densidad es continua.
  \item No tiene saltos puntuales.
  \item La pendiente se relaciona con la densidad.
\end{itemize}
\column{0.48\textwidth}
\begin{tikzpicture}
\begin{axis}[width=5cm,height=3.6cm, ymin=0,ymax=1, xmin=-3,xmax=3, xlabel={$x$}, ylabel={$F(x)$}]
\addplot[domain=-3:3, samples=100, blue, thick] {1/(1+exp(-1.7*x))};
\end{axis}
\end{tikzpicture}
\end{columns}
\end{frame}

\begin{frame}{Propiedades de la FDA}
\small
\begin{tabular}{p{0.35\textwidth}p{0.50\textwidth}}
\toprule
Propiedad & Interpretación \\
\midrule
Monotonía & si $a<b$, entonces $F(a)\le F(b)$ \\
Límite inferior & $\lim_{x\to-\infty}F(x)=0$ \\
Límite superior & $\lim_{x\to\infty}F(x)=1$ \\
Intervalos & $P(a<X\le b)=F(b)-F(a)$ \\
Discreta & los saltos recuperan $f(x)$ \\
Continua & si es derivable, $F'(x)=\rho(x)$ \\
\bottomrule
\end{tabular}
\end{frame}

\begin{frame}{Ejemplo con FDA}
\begin{block}{Dato}
Para una variable discreta, $F(1)=0{,}40$ y $F(3)=0{,}85$.
\end{block}
\pause
\[
P(1<X\le3)=F(3)-F(1)=0{,}85-0{,}40=0{,}45.
\]
\begin{exampleblock}{Interpretación}
La probabilidad acumulada entre 1 y 3 es $45\%$.
\end{exampleblock}
\end{frame}

\section{Esperanza matemática}

\begin{frame}{Esperanza matemática}
\begin{block}{Idea intuitiva}
La esperanza es el promedio a largo plazo de una variable aleatoria.
\end{block}
\begin{columns}[T]
\column{0.48\textwidth}
\begin{itemize}
  \item Centro de gravedad de la distribución.
  \item Valor medio esperado si el experimento se repite muchas veces.
  \item No siempre es un valor observable.
\end{itemize}
\column{0.48\textwidth}
\begin{tikzpicture}
\draw[thick] (0,0)--(4,0);
\fill[blue!20] (1.2,0) circle (.35);
\fill[blue!20] (2.0,0) circle (.35);
\fill[blue!20] (3.1,0) circle (.35);
\draw[red, very thick] (2.15,-.5)--(2.15,.5) node[above] {$E(X)$};
\end{tikzpicture}
\end{columns}
\end{frame}

\begin{frame}{Esperanza discreta}
\begin{block}{Fórmula}
\[
E(X)=\sum_x x\,f(x)
\]
\end{block}
\begin{itemize}
  \item $x$: valor que puede tomar la variable.
  \item $f(x)$: probabilidad de ese valor.
  \item $x f(x)$: aporte ponderado al promedio.
\end{itemize}
\end{frame}

\begin{frame}{Ejemplo: esperanza en un dado}
\small
Para un dado justo, $X=$ resultado.
\begin{align*}
E(X)&=1\cdot\frac16+2\cdot\frac16+3\cdot\frac16+4\cdot\frac16+5\cdot\frac16+6\cdot\frac16\\
&=\frac{21}{6}=3{,}5.
\end{align*}
\begin{alertblock}{Cuidado}
Aunque $3{,}5$ no sale en un dado, sí es el promedio esperado a largo plazo.
\end{alertblock}
\end{frame}

\begin{frame}{Esperanza continua}
\begin{block}{De suma a integral}
\[
E(X)=\int_{-\infty}^{\infty}x\rho(x)\,dx
\]
\end{block}
\begin{itemize}
  \item La integral suma infinitos aportes pequeños.
  \item La densidad actúa como ponderación local.
  \item Es el centro de gravedad de la curva.
\end{itemize}
\end{frame}

\begin{frame}{Ejemplo continuo: uniforme $U(0,10)$}
\[
E(X)=\int_0^{10}x\frac{1}{10}\,dx
=\frac{1}{10}\left[\frac{x^2}{2}\right]_0^{10}=5.
\]
\begin{exampleblock}{Interpretación}
Si la medición es uniforme entre 0 y 10, su centro esperado es 5.
\end{exampleblock}
\end{frame}

\begin{frame}{Esperanza de funciones}
\begin{block}{Definición}
Si interesa una transformación $g(X)$, entonces:
\[
E[g(X)]=\sum_x g(x)f(x)\quad\text{(discreta)}
\]
\[
E[g(X)]=\int g(x)\rho(x)\,dx\quad\text{(continua)}
\]
\end{block}
\begin{exampleblock}{Motivación}
Costos, pérdidas, energía, potencia o utilidad suelen ser funciones de una variable aleatoria.
\end{exampleblock}
\end{frame}

\begin{frame}{Propiedades de la esperanza}
\begin{theorem}[Linealidad]
\[
E(aX+b)=aE(X)+b
\]
\[
E(aX+bY)=aE(X)+bE(Y)
\]
\end{theorem}
\begin{block}{Interpretación}
Cambiar escala multiplica el promedio; desplazar la variable desplaza el promedio.
\end{block}
\end{frame}

\begin{frame}{Mediana, moda y esperanza}
\small
\begin{tabular}{p{0.25\textwidth}p{0.35\textwidth}p{0.28\textwidth}}
\toprule
Medida & Definición probabilística & Sensibilidad \\
\midrule
Esperanza & promedio ponderado & sensible a colas \\
Mediana & deja 50\% a cada lado & robusta \\
Moda & valor más probable & puede no ser única \\
\bottomrule
\end{tabular}
\begin{alertblock}{Error común}
Creer que esperanza, mediana y moda siempre coinciden. Coinciden en distribuciones muy simétricas, pero no en general.
\end{alertblock}
\end{frame}

\section{Varianza}

\begin{frame}{¿Por qué la esperanza no basta?}
\begin{columns}[T]
\column{0.48\textwidth}
\begin{block}{Dos variables pueden tener igual esperanza}
pero distinta dispersión alrededor de ese centro.
\end{block}
\column{0.48\textwidth}
\begin{tikzpicture}
\begin{axis}[width=5cm,height=3.6cm, ymin=0,ymax=.55, xmin=-4,xmax=4, xlabel={$x$}, ylabel={}]
\addplot[domain=-4:4, samples=80, blue, thick] {1/sqrt(2*pi)*exp(-x^2/2)};
\addplot[domain=-4:4, samples=80, red, thick, dashed] {1/sqrt(2*pi*4)*exp(-x^2/8)};
\end{axis}
\end{tikzpicture}
\end{columns}
\end{frame}

\begin{frame}{Varianza: definición}
\begin{block}{Definición formal}
\[
\operatorname{Var}(X)=E[(X-\mu)^2],\qquad \mu=E(X).
\]
\end{block}
\begin{itemize}
  \item Mide dispersión probabilística respecto de la esperanza.
  \item Usa cuadrados para evitar cancelaciones.
  \item Queda en unidades al cuadrado.
\end{itemize}
\end{frame}

\begin{frame}{Forma equivalente de la varianza}
\begin{theorem}[Fórmula práctica]
\[
\operatorname{Var}(X)=E(X^2)-\mu^2
\]
\end{theorem}
\begin{block}{Demostración intuitiva}
Expandir $(X-\mu)^2=X^2-2\mu X+\mu^2$ y aplicar linealidad de la esperanza.
\end{block}
\end{frame}

\begin{frame}{Ejemplo: varianza de un dado}
\small
Para un dado justo, $\mu=3{,}5$.
\[
E(X^2)=\frac{1^2+2^2+3^2+4^2+5^2+6^2}{6}=\frac{91}{6}
\]
\[
\operatorname{Var}(X)=\frac{91}{6}-(3{,}5)^2=2{,}9167.
\]
\begin{exampleblock}{Lectura}
La dispersión típica se obtiene con la desviación estándar.
\end{exampleblock}
\end{frame}

\begin{frame}{Desviación estándar}
\begin{block}{Definición}
\[
\sigma=\sqrt{\operatorname{Var}(X)}
\]
\end{block}
\begin{itemize}
  \item Vuelve a las unidades originales.
  \item Facilita interpretar la dispersión.
  \item Para el dado: $\sigma\approx\sqrt{2{,}9167}=1{,}71$.
\end{itemize}
\end{frame}

\begin{frame}{Propiedades de la varianza}
\begin{theorem}[Cambio de escala]
\[
\operatorname{Var}(aX+b)=a^2\operatorname{Var}(X)
\]
\end{theorem}
\begin{itemize}
  \item Sumar $b$ desplaza la variable, pero no cambia su dispersión.
  \item Multiplicar por $a$ escala las distancias por $a$.
  \item La varianza escala por $a^2$ porque mide cuadrados.
\end{itemize}
\end{frame}

\section{Síntesis}

\begin{frame}{Mapa conceptual completo}
\centering
\begin{tikzpicture}[node distance=.58cm, box/.style={rectangle, rounded corners, draw=blue!65, fill=blue!8, align=center, minimum width=3.1cm, minimum height=.52cm}, arrow/.style={-{Latex}, thick}]
\node[box] (a) {Experimento aleatorio};
\node[box, below=of a] (b) {Espacio muestral $\Omega$};
\node[box, below=of b] (c) {Variable aleatoria $X$};
\node[box, below=of c] (d) {Distribución};
\node[box, below=of d] (e) {Probabilidad};
\node[box, below=of e] (f) {Distribución acumulada $F$};
\node[box, below=of f] (g) {Esperanza $E(X)$};
\node[box, below=of g] (h) {Varianza $\operatorname{Var}(X)$};
\draw[arrow] (a)--(b); \draw[arrow] (b)--(c); \draw[arrow] (c)--(d); \draw[arrow] (d)--(e); \draw[arrow] (e)--(f); \draw[arrow] (f)--(g); \draw[arrow] (g)--(h);
\end{tikzpicture}
\end{frame}

\begin{frame}{Tabla comparativa final}
\tiny
\begin{tabular}{p{0.27\textwidth}p{0.32\textwidth}p{0.32\textwidth}}
\toprule
Comparación & Primera idea & Segunda idea \\
\midrule
Discreta vs continua & conteos, puntos & mediciones, intervalos \\
Probabilidad vs densidad & $f(x)=P(X=x)$ & $\rho(x)$ genera áreas \\
Puntual vs acumulada & valor exacto & $F(x)=P(X\le x)$ \\
Esperanza vs media descriptiva & promedio teórico & promedio observado \\
Varianza probabilística vs descriptiva & dispersión del modelo & dispersión de datos \\
\bottomrule
\end{tabular}
\end{frame}

\begin{frame}{Errores frecuentes}
\small
\begin{enumerate}
  \item Confundir resultado $\omega$ con valor $X(\omega)$.
  \item Confundir dominio $\Omega$ con rango $R_X$.
  \item Tratar una densidad como probabilidad puntual.
  \item Olvidar que en variables continuas $P(X=x)=0$.
  \item Calcular $P(a<X\le b)$ sin usar $F(b)-F(a)$.
  \item Interpretar esperanza como resultado garantizado.
  \item Comparar varianzas sin revisar unidades.
\end{enumerate}
\end{frame}

\begin{frame}{Formulario completo de la unidad}
\small
\begin{columns}[T]
\column{0.48\textwidth}
\begin{align*}
X&:\Omega\to\mathbb{R}\\
f(x)&=P(X=x)\\
\sum_x f(x)&=1\\
F(x)&=P(X\le x)\\
P(a<X\le b)&=F(b)-F(a)
\end{align*}
\column{0.48\textwidth}
\begin{align*}
P(a<X<b)&=\int_a^b\rho(x)dx\\
E(X)&=\sum_x xf(x)\\
E(X)&=\int x\rho(x)dx\\
\operatorname{Var}(X)&=E(X^2)-\mu^2\\
\sigma&=\sqrt{\operatorname{Var}(X)}
\end{align*}
\end{columns}
\end{frame}

\begin{frame}{Ejercicios integradores}
\small
\begin{enumerate}
  \item Lanzar dos monedas. Defina $X=$ número de caras, construya $R_X$ y $f(x)$.
  \item Para un dado justo, calcule $P(X\le4)$ usando la FDA.
  \item Si $X\sim U(0,10)$, calcule $P(3<X<7)$ y explique el área.
  \item Un sensor registra número de alarmas por hora. ¿La variable es discreta o continua?
  \item Compare dos sistemas con igual esperanza de falla pero distinta varianza.
\end{enumerate}
\end{frame}

\begin{frame}{Preguntas conceptuales de cierre}
\begin{enumerate}
  \item ¿Por qué una variable aleatoria es una función?
  \item ¿Qué se pierde si solo se informa la esperanza?
  \item ¿Cuándo conviene usar una FDA en vez de una función puntual?
  \item ¿Por qué la densidad puede superar 1 sin violar axiomas?
  \item ¿Cómo conecta esta unidad con inferencia estadística?
\end{enumerate}
\end{frame}

\begin{frame}{Cierre de variables aleatorias}
\centering
\Large
Una variable aleatoria traduce incertidumbre en matemáticas operables.
\vspace{0.8em}

\normalsize
Distribución, acumulada, esperanza y varianza son el puente hacia modelos probabilísticos e inferencia.
\end{frame}

\end{document}
